\documentclass{article}
\usepackage{iclr2027_conference,times}


\usepackage{amsmath,amssymb,amsfonts}
\usepackage{graphicx}
\usepackage{textcomp}
\usepackage{xcolor}
\usepackage{booktabs}
\usepackage{array}
\usepackage{multirow}
\usepackage{needspace}
\usepackage{algorithm}
\usepackage[noend]{algpseudocode}
\algrenewcommand\algorithmicindent{1em}
\usepackage{wrapfig}
\usepackage[capitalise,noabbrev]{cleveref}
\crefname{figure}{Fig.}{Figs.}
\Crefname{figure}{Fig.}{Figs.}



\newcommand{\hdr}[2]{\begin{tabular}[b]{@{}c@{}}\textbf{#1}\\\textbf{#2}\end{tabular}}
\newcommand{\thead}[1]{\textbf{\shortstack[b]{#1}}}
\usepackage[hyphens]{url}


\usepackage{xurl}

\raggedbottom


\begin{document}
\title{LaDiM: Layered Diagnosis for Multi-Agent Code Migration}
\author{Anonymous Authors}
\maketitle
\suppressfloats[t]
\begin{abstract}
Migrating code across deep learning frameworks enables the reuse of training programs in new software and hardware environments. Different frameworks, however, might compute gradients or update parameters in different ways, leading to wrong results even when a translated program runs without error. Current methods are mostly designed for language transfer, neglecting the subtle differences in framework level. To address this, LaDiM preserves the underlying training computation order, e.g., forward, gradient computation, update, and uses the dependencies among these stages to detect and localize silent computational errors. LaDiM includes Translator, Verifier, and Repair agents working under a shared Orchestrator. The Translator produces an initial output. The Verifier traces discrepancies along the computation order and hands its findings to the Repair Agent, which revises the candidate using this evidence and a persistent repair history. For repository-scale migrations, LaDiM further introduces repository context management which tracks cross-file dependencies and preserves diagnostic context as agents move across the repository. Evaluation across framework migration, language migration, and repository migration tasks shows that LaDiM achieves comparable or better migration results with up to 80.8\% fewer tokens than state-of-the-art methods, demonstrating its effectiveness.
\end{abstract}

\section{Introduction}
{\hyphenpenalty=10000\emergencystretch=1em
\begin{wrapfigure}{r}{0.43\textwidth}
\vspace{-\intextsep}
\centering
\includegraphics[width=\linewidth]{figures/migration_motivation.pdf}
\caption{Tokens used on 50 PyTorch to MindSpore migration tasks.}
\label{fig:migration-motivation}
\end{wrapfigure}
Code development environments and languages have grown increasingly diverse. Focusing on deep learning, popular frameworks such as PyTorch, TensorFlow, and MindSpore provide programming interfaces, execution semantics, and hardware backends~\citep{paszke2019pytorch,abadi2016tensorflow,mindspore2020whitepaper}. To enable easy adaptation of deep models and training pipelines on new software and hardware environments, migrating existing code across languages and frameworks can reduce development cost~\citep{guan2025repotransagent}. This need has driven the development of numerous code translation methods. Large language models (LLMs) reduce the manual effort involved in code translation~\citep{he-etal-2025-execoder,macedo2024intertrans}, while automated tests and interactive repair help assess and correct translated behavior~\citep{roziere2021unittests,yang2024sweagent,ibrahimzada2025matchfixagent}.

However, success at translating code does not guarantee equal translation performance across deep learning frameworks, for example, migrating from PyTorch to MindSpore. Different frameworks use different interfaces and execution mechanisms for forward computation, gradient computation, and parameter updates~\citep{paszke2019pytorch,mindspore2024gradient}. Further differences in operator semantics, automatic differentiation, and optimizer interfaces can change model behavior after migration~\citep{paszke2019pytorch,mindspore2020whitepaper}. As a result, a translated program can run without raising error while still producing incorrect results due to silent failures in computing losses, gradients or parameter updates. We use semantic drift to describe the changes in such training behavior introduced by migration. These methods primarily target language level translation and test based functional correctness, leaving framework specific computational drift less directly addressed~\citep{he-etal-2025-execoder,macedo2024intertrans,ibrahimzada2025matchfixagent}.

To address this, we aim to locate and repair the true cause of semantic drift after initial translation. We propose LaDiM (Layered Diagnosis for Multi-Agent Code Migration), a multi-agent framework that preserves the training computation order of deep learning programs. We propose Layered Diagnosis, inspecting the translated candidate stage by stage, e.g., forward computation, gradient computation, and parameter update, in the order of how training actually executes. The observed discrepancy during inspection is handed to a separate Repair Agent and saved to local repair history for future use in later rounds. For repository-scale migration, we propose repository context management that first organizes code files into a dependency graph and repairs them in topological order to avoid redundant inspection and improve effectiveness. An Orchestrator schedules overall process and returns verification results to guide further repair.

LaDiM achieves comparable or better translation results with substantially lower tokens consumed. As shown in \cref{fig:migration-motivation}, LaDiM completes all 50 migrations from PyTorch to MindSpore with 57.4\% fewer total tokens than MatchFixAgent, and 75.2\% fewer than SWE-agent while obtaining higher accuracy. On the repository-level translation tasks, LaDiM achieves the best translation accuracy while saving more than 80\% tokens used. Our contributions are as follows:
\par}
\begin{itemize}
\item We develop Layered Diagnosis, which uses dependencies among execution, forward values, gradients, and parameter updates to locate the causes of changed training behavior.
\item We design an agent framework with independent evidence handoff and repair history to connect investigation, edits, and verification. Repository context management coordinates this process across dependent files through Repository Structural Analysis and Repair Dependency Graph Planning.
\item We demonstrate LaDiM's effectiveness and token efficiency across framework, language, and repository migration, and identify the contributions of training signals and agent components through ablation studies.
\end{itemize}

\section{Related Work}

\paragraph{Code Translation and Framework Migration.}
Neural translation methods learn code mappings from multilingual corpora and use automated tests or compiler representations to improve translations~\citep{roziere2020transcoder,roziere2021unittests,szafraniec2023compiler}. ExeCoder encodes functional semantics, syntax structure, and variable dependencies~\citep{he-etal-2025-execoder}, while InterTrans explores intermediate translation paths~\citep{macedo2024intertrans}. MindConverter converts model definitions and weights and maps supported APIs~\citep{mindspore2023mindconverter}. RepoTransAgent uses repository context, test failures, and reflection to refine translated code~\citep{guan2025repotransagent}. LaDiM investigates how translated training code computes forward values and gradients and updates model parameters. At repository scale, it coordinates changes to shared interfaces and checks their effects on the repository entry points that use them.

\paragraph{Differential Testing and Deep Learning Verification.}
Differential testing exposes compiler and program faults through disagreement between related executions~\citep{yang2011csmith,le2014emi}. DeepXplore detects behavioral differences between deep learning systems~\citep{pei2017deepxplore}, while DeepGauge measures test coverage in these systems~\citep{ma2018deepgauge}. More directly related to framework migration, TensorScope tests corresponding APIs across frameworks to identify inconsistent behavior~\citep{deng2023tensorscope}. NablaFuzz compares execution scenarios to test automatic differentiation, including first-order and higher-order gradients~\citep{yang2023nablafuzz}. LaDiM uses training differences as evidence for code repair. Its Verifier Agent inspects code to explain discrepancies in forward values, gradients, and parameter updates, then passes the supporting observations and hypotheses to the Repair Agent for testing during repair.

\paragraph{Automated Program Repair and LLM Agents.}
Automated program repair uses tests or semantic constraints to guide patch generation~\citep{nguyen2013semfix}. Pretrained language models generate patches from buggy code and its context, with tests used to evaluate candidate patches~\citep{xia2023aprplm}. SWE-agent supports repository navigation, editing, and testing~\citep{yang2024sweagent}. Reflexion and Self-Debugging use execution feedback to support reflection and self-correction~\citep{shinn2023reflexion,chen2023selfdebugging}. CodeAct uses executable code actions and their results to guide subsequent actions~\citep{wang2024codeact}. ChatDev and MetaGPT assign software development activities to specialized roles~\citep{qian2023chatdev,hong2023metagpt}. AutoGen and AgentVerse provide general frameworks for multi-agent collaboration~\citep{wu2023autogen,chen2023agentverse}. MatchFixAgent takes a source program and an existing translation and uses several semantic analyses to guide test generation and repair~\citep{ibrahimzada2025matchfixagent}. LaDiM includes initial translation and organizes diagnosis around dependencies in training computations. The Verifier Agent passes its findings and supporting observations to a separate Repair Agent, which uses measurements from each submitted revision to revise its hypotheses.

\section{Methods}
\label{sec:method}
\begin{figure}[!t]
\centering
\includegraphics[width=0.9\textwidth]{figures/hierarchical_feedback_architecture.pdf}
\caption{Overview of LaDiM. (a) The Translator produces a candidate, and the Orchestrator compares its execution, forward values, gradients, and parameter updates with the source. The Verifier investigates the first divergent stage (gradients here) and hands code observations, test results, locations, and hypotheses to the Repair Agent. The Repair Agent tests and edits the candidate, using the Orchestrator's feedback until acceptance or budget exhaustion. (b) Repository context management analyzes code structure, orders repairs through work-unit dependencies, and reconstructs compact contexts from archived evidence.}
\label{fig:architecture}
\end{figure}

\subsection{Preliminary}
A migration task $\mathcal T=(P,L_s,F_s,L_t,F_t,W)$ specifies a source program $P$, its language and framework $(L_s,F_s)$, the requested target environment $(L_t,F_t)$, and an editable file scope $W$. Migration produces a target program $\hat P$ that preserves the source's training behavior. The scope can include the translated program, its supporting library, or a repository. An edit modifies code within $W$.

At step $t$, we record source observations $(e_t,v_t,g_t,\theta_{t+1}-\theta_t)$ and target observations $(\hat e_t,\hat v_t,\hat g_t,\hat\theta_{t+1}-\hat\theta_t)$, with numerical values recorded when available. Here $e_t$ records the execution outcome, including success or an exception, $v_t$ contains selected forward values and losses, $g_t$ contains gradients, and $\theta_t$ denotes the parameters before the update. After aligning corresponding tensors and parameters, we define
\begin{align}
d_{\mathrm{execution},t}&=\mathbf{1}[e_t\ne\hat e_t], &
d_{\mathrm{forward},t}&=v_t-\hat v_t,\\
d_{\mathrm{gradient},t}&=g_t-\hat g_t, &
d_{\mathrm{update},t}&=(\theta_{t+1}-\theta_t)-(\hat\theta_{t+1}-\hat\theta_t).
\end{align}
For a numerical comparison $d=x-\hat x$, $d\simeq0$ denotes agreement within its tolerance, such as $|d_j|\leq\alpha+\beta|x_j|$ for absolute tolerance $\alpha$ and relative tolerance $\beta$, or a specified norm bound. Forward values determine the loss, differentiation produces the gradients, and the optimizer uses those gradients and its state to update parameters. The updated parameters then affect subsequent forward computations.

The LLM $M$ performs translation, investigation, and repair under role instructions $A_T$, $A_V$, and $A_R$. The Orchestrator $\mathcal O$ evaluates candidates and manages their shared investigation and repair budget $B$. In the algorithms, $S.\hat P$ is the candidate, $S.m$ contains its measurements, $q$ is its verification status, and $H$ is a conversation. Context $C$ retains evidence and provides tools; for repository tasks, it also stores the dependency plan and local check results. A submission sends the complete candidate for external verification after a repair stage that can contain several LLM calls, tests, and edits. Each $\textsc{LLMCall}(M,A,H,b,B)$ invokes $M$ with instructions $A$ and conversation $H$, charging its usage to the stage allowance $b$ and shared budget $B$. The Orchestrator counts each submission against $B$.

\subsection{Verifier Agent with Layered Diagnosis}
\label{sec:hierarchy}

{
\setlength{\intextsep}{0pt}
\begin{wrapfigure}[16]{l}{0.50\textwidth}
\begin{minipage}{\linewidth}
\begin{algorithm}[H]
\caption{Layered Diagnosis}
\label{alg:diagnose}
\raggedright
\small
\begin{algorithmic}[1]
\Require $\mathcal T,S,M,B,C$
\Ensure Evidence $E$
\Function{LayeredDiagnosis}{$\mathcal T,S,M,B,C$}
  \State $H\gets\Call{StartConversation}{A_V,\mathcal T,S}$
  \State $b\gets\Call{DiagnosisAllowance}{B}$
  \State $\mathcal D:\ \mathrm{execution}\to v_t\to g_t\to\Delta\theta_t\to v_{t+1}$
  \While{$\Call{WithinLimits}{b,B}$}
    \State $H\gets C.\Call{PrepareContext}{H,S}$
    \State $(\ell,h,a)\gets\Call{LLMCall}{M,A_V,H,b,B}$
    \State $H\gets\Call{ToolFeedback}{a,H,S,A_V,C}$
    \State \textbf{if} $a.\mathrm{complete}$ \textbf{then break}
  \EndWhile
  \State \Return $\Call{PublicEvidence}{H,C}$
\EndFunction
\Function{ToolFeedback}{$a,H,S,A,C$}
  \State $(C,S,r)\gets C.\Call{ExecuteTools}{a,S,A}$
  \State \Return $H\mathbin{\|}(a,r)$
\EndFunction
\end{algorithmic}
\end{algorithm}
\end{minipage}
\end{wrapfigure}
The Verifier Agent receives the Translator Agent's candidate and the Orchestrator's measurements (\cref{fig:architecture}a). It uses Layered Diagnosis to focus its investigation on the computations responsible for observed discrepancies. For an \textbf{execution discrepancy} ($d_{\mathrm{execution},t}=1$), it inspects the exception and the calls that produced it. For a \textbf{forward discrepancy} ($d_{\mathrm{forward},t}\not\simeq0$), it traces the operators that produce different values and tests the suspected computation. When forward values agree but a \textbf{gradient discrepancy} ($d_{\mathrm{gradient},t}\not\simeq0$) remains, it examines parameter registration and the path from the parameters to the loss. When both forward values and gradients agree, it investigates an \textbf{update discrepancy} ($d_{\mathrm{update},t}\not\simeq0$) through the optimizer's rule and state. The first step with an observed discrepancy further guides this investigation because a changed update can alter later forward computations. The LLM selects code inspections and tests that distinguish possible causes, then records the code locations and observations supporting its hypotheses. These records form the evidence for the independent handoff to the Repair Agent.

\cref{alg:diagnose} expresses this reasoning through the training dependency relation $\mathcal D$, where $\Delta\theta_t=\theta_{t+1}-\theta_t$. The symbols $\ell$ and $h$ summarize the LLM's choice of computation and hypothesized cause within a call, and $a$ is its tool request or completion response. The Verifier uses both agreement and disagreement between layers to select $\ell$, then tests $h$ against code and execution evidence. The shared \textsc{ToolFeedback} operation executes requests in order and adds their results to the conversation. During investigation, the Verifier Agent can create and run scratch tests while candidate and supporting library code remain unchanged.\par}

\subsection{Repair Agent with Evidence Handoff and History}
\label{sec:repair-agent}

{
\setlength{\intextsep}{0pt}

Independent evidence handoff gives the Repair Agent the Verifier Agent's recorded code observations, test results, and hypotheses, identified as the Verifier Agent's findings. The Repair Agent reassesses this evidence, runs further tests, and edits files in $W$, including supporting library code when the task permits it. The Orchestrator verifies each submission and returns the remaining discrepancies to guide the Repair Agent's next investigation and edits.
\Needspace{13\baselineskip}
\begin{wrapfigure}[13]{l}{0.55\textwidth}
\begin{minipage}{\linewidth}
\begin{algorithm}[H]
\caption{\textsc{Repair}$(\mathcal T,S,E,q,M,B,C)$}
\label{alg:repair}
\raggedright
\small
\begin{algorithmic}[1]
\State $H\gets\Call{IndependentHandoff}{A_R,\mathcal T,E,C}$
\While{$\Call{SubmissionsLeft}{B}\land\allowbreak\Call{BudgetLeft}{B}$}
  \State $H\gets H\mathbin{\|}S.m$; $b\gets\Call{RepairAllowance}{B}$
  \While{$\Call{WithinLimits}{b,B}$}
    \State $H\gets C.\Call{PrepareContext}{H,S}$
    \State $a\gets\Call{LLMCall}{M,A_R,H,b,B}$ \Comment{Test and edit}
    \State $H\gets\Call{ToolFeedback}{a,H,S,A_R,C}$
    \State \textbf{if} $\Call{RepairComplete}{a,H}$ \textbf{then break}
  \EndWhile
  \State $(S.m,q)\gets\mathcal O.\Call{Submit}{S.\hat P,B}$
  \State \textbf{if} $q=\mathrm{accepted}$ \textbf{then return} $(S.\hat P,q)$
\EndWhile
\State \Return $(S.\hat P,q)$
\end{algorithmic}
\end{algorithm}
\end{minipage}
\end{wrapfigure}
Repair history preserves earlier edits, observations, and hypotheses across attempts. The continuing conversation includes failed tests and verification feedback, allowing the agent to relate a changed result to the preceding modification and revise explanations that did not resolve the fault. The current candidate also persists, so later attempts build on applied changes.

\cref{alg:repair} reuses \textsc{ToolFeedback} from \cref{alg:diagnose}, with $H\mathbin{\|}x$ denoting the addition of a record $x$ to the repair history. Its \textsc{RepairComplete} condition requires a completion response, an applied edit in the current stage, and a test executed after the last edit. If an edit or test is missing, the conversation requests that action within the remaining budget. The Orchestrator evaluates the candidate when the repair stage ends.\par}

\begin{algorithm}[!htb]
\caption{Repository migration and repair coordination}
\label{alg:repository}
\small
\begin{algorithmic}[1]
\Function{RepositoryMigrate}{$\mathcal T,M,B$}
  \State $C\gets\Call{RepositoryContext}{\mathcal T}$
  \State $S.\hat P\gets\Call{LLMTranslate}{M,A_T,\mathcal T}$
  \State $(S.m,q)\gets\mathcal O.\Call{Verify}{S.\hat P}$ \Comment{Complete repository}
  \State \textbf{if} $q=\mathrm{accepted}$ \textbf{then return} $(S.\hat P,q)$
  \State $E\gets\Call{LayeredDiagnosis}{\mathcal T,S,M,B,C}$ \Comment{\cref{alg:diagnose}}
  \State \Return $\Call{Repair}{\mathcal T,S,E,q,M,B,C}$ \Comment{\cref{alg:repair}}
\EndFunction
\Function{$C$.ExecuteTools}{$a,S,A$}
  \State $r\gets[\,]$
  \ForAll{$c\in a.\mathrm{tools}$}
    \State $U\gets\varnothing$
    \State \textbf{if} $c.\mathrm{kind}=\mathrm{structure}$ \textbf{then} $o\gets\Call{RepositoryStructuralAnalysis}{C,S}$
    \State \textbf{else if} $c.\mathrm{kind}=\mathrm{plan}$ \textbf{then} $(C.G,U,o)\gets\Call{UpdateDependencyGraph}{c,C}$
    \State \textbf{else} $(C,S,o,U)\gets\Call{Execute}{c,C,S,A}$
    \State $C.E\gets C.E\mathbin{\|}(o,\Call{CodeVersion}{S})$; $r\gets r\mathbin{\|}o$
    \State $C.Q[U\cup\Call{Dependents}{C.G,U}]\gets\mathrm{invalid}$
  \EndFor
  \State \Return $(C,S,r)$
\EndFunction
\end{algorithmic}
\end{algorithm}

\subsection{Repository Coordination}
\label{sec:repository-method}
Repository migration raises two problems that single-program migration does not: a repair may span several interdependent files, and evidence gathered on one file must remain available while the agent works on others. Repository context management addresses both through three components (\cref{fig:architecture}b).

\textbf{Repository Structural Analysis.} This tool returns an inventory of the target repository: its files, the imports among them, the functions and classes each file defines, and the cells of each notebook. It also lists files not yet assigned to any work unit. The agent invokes it on demand.

\textbf{Repair Dependency Graph Planning.} Based on this inventory, the agent partitions the target files into \emph{work units}. Each unit specifies its editable files, a repair goal, its tests, and the units it depends on; together, the units form a directed acyclic graph $C.G$. A unit may start only when the checks of its prerequisites remain valid, so shared implementations are repaired before the units that use them. $C.Q$ stores each unit's syntax check and test results with the code version on which they ran; when a file or the plan changes, the results of the affected units and their dependents are invalidated. Unit checks only guide the agent; the Orchestrator decides acceptance by verifying the complete repository.

\textbf{Evidence archive and context reconstruction.} Tool observations, code read by the agent, measurements, hypotheses, and conversations are stored in an archive $C.E$ outside the editable workspace, each tagged with the code version it describes. At evidence handoff, when the agent switches work units, or when the conversation exceeds 75\% of the context capacity, \textsc{PrepareContext} replaces the conversation with a compact one containing the current plan and progress, the structured findings, the latest measurements, and previously read code that is relevant to the active unit and unchanged since it was read.

\cref{alg:repository} connects these components to \cref{alg:diagnose,alg:repair}. \textsc{Execute} handles unit selection, file and notebook-cell operations, tests, and evidence retrieval, and returns an observation $o$ with the set $U$ of units affected by file or plan changes. All units share the repository files, the repair history, and the total budget.



\section{Experiments}

\subsection{Experimental Setup}
\label{sec:setup}

\textbf{Tasks.} The main comparison contains 50 migration tasks from PyTorch to MindSpore, covering tensor operations, CNNs, MLPs, Transformers, language models, and parameter and optimizer interfaces. We also migrate 18 Java/DJL programs to Python/PyTorch, comprising eight textbook examples and ten applications~\citep{djl2026,d2ljava2026}, and two public repositories for time series~\citep{jingles2026timeseries} and recommendation~\citep{chak2026twotower}. The repositories contain 13 and 24 code files or notebooks, totaling 865 and 2,086 nonblank lines.

To examine repair of actual translation errors, we use ten programs produced by the initial translation, of which five fail evaluation and five pass before repair. We use 16 faulty candidates to study training signals and six faulty MLP and CNN candidates to study repair on JAX. A separate detection study examines how early gradient, update, and loss checks reveal training faults. \cref{sec:task-construction} summarizes the program collections in \cref{tab:collections}.

\textbf{Baselines.} We compare LaDiM with SWE-agent~\citep{yang2024sweagent} and MatchFixAgent~\citep{ibrahimzada2025matchfixagent} using a common DeepSeek backend and each baseline's native repair procedure. The main table also includes Direct LLM for initial translation, CodeTransEngine's direct translation configuration~\citep{macedo2025codetransengine}, MSAdapter~\citep{openi2025msadapter}, and InterTrans~\citep{macedo2024intertrans}. Test-guided repair rewrites the complete candidate using test feedback. Direct repair in the analysis studies receives the same candidate as LaDiM and uses LaDiM's inspection, testing, and editing tools. For JAX, Ivy and \texttt{torch2jax} each perform one native conversion of the faulty candidates (\cref{sec:jax-protocol}).

\textbf{Evaluation Protocol.} Methods use common source programs, evaluation inputs, and corresponding initial parameters. Repair methods also receive the same initial translations. Acceptance requires agreement in execution, forward values, gradients, and parameter updates, together with the task's optimizer and interface checks. For repositories, \emph{Checks} counts successful execution, numerical, and interface checks. \emph{Entries} counts successful execution of the training script, notebook, and teaching examples. \emph{Tests} counts passing tests supplied by the repository. Candidates that pass the evaluation available during repair are checked on additional seeds. Investigation and repair allow up to 40 LLM calls per program and 80 per repository, with at most four submissions. End-to-end costs include initial translation once and all subsequent calls, including failed attempts. \cref{sec:supplementary-protocols} gives the numerical tolerances, seeds, and complete budgets.

\subsection{Main Results}
\label{sec:main-results}
\begin{table}[t]
\caption{Migration acceptance and token use on common inputs: (a,b) individual programs; (c,d) repositories. Tokens are in millions. Calls count all LLM calls in (a,b) and investigation and repair calls in (c,d).}
\label{tab:main-results}
\centering
\small
\setlength{\tabcolsep}{3pt}

\begin{minipage}[t]{0.485\linewidth}
\vspace{0pt}
\raggedright\textit{(a) PyTorch to MindSpore}\par\smallskip
\setlength{\tabcolsep}{2pt}
\begin{tabular*}{\linewidth}{@{\extracolsep{\fill}}>{\raggedright\arraybackslash}p{2.8cm}rrr@{}}
\toprule
\textbf{Method} & \textbf{Calls} & \textbf{Tokens} & \textbf{Accepted} \\
\midrule
Direct LLM & 29 & 0.429 & 29/50 \\
CodeTransEngine & 29 & 0.280 & 31/50 \\
MSAdapter & -- & -- & 15/50 \\
SWE-agent & 1,097 & 20.831 & 44/50 \\
MatchFixAgent & 619 & 12.097 & \textbf{50/50} \\
\textbf{LaDiM (ours)} & 326 & 5.159 & \textbf{50/50} \\
\bottomrule
\end{tabular*}
\end{minipage}
\hfill
\begin{minipage}[t]{0.485\linewidth}
\vspace{0pt}
\raggedright\textit{(b) Java/DJL to Python/PyTorch}\par\smallskip
\setlength{\tabcolsep}{2pt}
\begin{tabular*}{\linewidth}{@{\extracolsep{\fill}}>{\raggedright\arraybackslash}p{2.8cm}rrr@{}}
\toprule
\textbf{Method} & \textbf{Calls} & \textbf{Tokens} & \textbf{Accepted} \\
\midrule
Direct LLM & 18 & 0.694 & 1/18 \\
InterTrans & 89 & 2.286 & 1/18 \\
Test-guided repair & 71 & 4.475 & 7/18 \\
SWE-agent & 682 & 24.321 & 8/18 \\
MatchFixAgent & 524 & 28.355 & 8/18 \\
\textbf{LaDiM (ours)} & 345 & 21.479 & \textbf{9/18} \\
\bottomrule
\end{tabular*}
\end{minipage}
\par\medskip
\begin{tabular*}{\linewidth}{@{\extracolsep{\fill}}lrrrrrrrr@{}}
\toprule
& \multicolumn{4}{c}{\textit{(c) Time series repository}} & \multicolumn{4}{c}{\textit{(d) Recommendation repository}} \\
\cmidrule(lr){2-5}\cmidrule(l){6-9}
\textbf{Method} & \textbf{Checks} & \textbf{Entries} & \textbf{Calls} & \textbf{Tokens} & \textbf{Checks} & \textbf{Tests} & \textbf{Calls} & \textbf{Tokens} \\
\midrule
SWE-agent & 69/69 & 3/3 & 80 & 8.200 & 104/145 & 10/10 & 80 & 8.317 \\
MatchFixAgent & 69/69 & 3/3 & 69 & 13.598 & 129/145 & 0/10 & 80 & 8.679 \\
\textbf{LaDiM (ours)} & 69/69 & 3/3 & \textbf{35} & 2.614 & \textbf{130/145} & 10/10 & 80 & 8.811 \\
\bottomrule
\end{tabular*}


\end{table}

\textbf{Migration to MindSpore.} LaDiM completes all 50 MindSpore migrations with 57.4\% fewer tokens than MatchFixAgent (\cref{tab:main-results}a). The savings cover both programs that require repair and programs that pass evaluation before repair (\cref{fig:repair-comparison}). Reduced input token usage accounts for 93.0\% of the savings, so most of the reduction comes from processing less context during investigation and repair. With submission limits of one, two, and four, LaDiM accepts 46/50, 50/50, and 50/50 tasks, respectively.
\begin{figure}[!t]
\centering
\includegraphics[width=0.96\linewidth]{figures/repair_comparison.pdf}
\caption{MindSpore migration token costs. (a) Initial translation and subsequent calls, grouped by whether repair is needed. (b) MatchFixAgent minus LaDiM tokens for 29 distinct inputs: tasks 1--20 pass before repair; tasks 21--29 require repair. Savings are sorted within each group.}
\label{fig:repair-comparison}
\end{figure}

\textbf{Migration across languages.} LaDiM accepts all programs completed by MatchFixAgent and an additional recurrent network example in the Java/DJL collection (\cref{tab:main-results}b). Its 9/18 accepted programs comprise all eight textbook examples and one application, compared with 8/18 for SWE-agent and MatchFixAgent. LaDiM uses 11.7\% fewer tokens than SWE-agent and 24.2\% fewer than MatchFixAgent.

\textbf{Repository migration.} All three agents complete the time series migration, preserving the shared model across a training script, notebook, and teaching examples. At the same complete acceptance of 69/69 checks, LaDiM uses 68.1\% fewer end-to-end tokens than SWE-agent and 80.8\% fewer than MatchFixAgent.

On recommendation, LaDiM restores reward model inference and retrieval, which remain unresolved in both baselines, and executes the training command. It achieves the highest behavioral coverage, passing 130/145 checks and all ten original tests. \cref{sec:repository-results} gives the remaining numerical differences and detailed checks.

\subsection{Analysis Experiments}
\label{sec:fault-characterization}

\textbf{Repairing Translation Errors.} Methods may edit both the ten translated programs and their supporting library. \cref{tab:repair-results}a reports how many failed programs are \emph{Repaired}, how many passing programs are \emph{Preserved}, and how many programs are finally \emph{Accepted}. LaDiM repairs four of the five failed programs and preserves all five programs that pass evaluation before repair, raising final acceptance from five to nine. Direct repair, SWE-agent, and MatchFixAgent retain the initial five successes but recover none of the failed programs. \cref{sec:natural-protocol} details a repaired LSTM whose recurrent operators and tensor indexing in the supporting library required two further submissions.
\begin{table}[t]
\caption{Repair of (a) ten initial translations, five failing and five passing, and (b) six faulty JAX candidates. LaDiM and Direct repair allow four submissions; Ivy and \texttt{torch2jax} perform one native conversion. Tokens are in millions for all investigation and repair calls.}
\label{tab:repair-results}
\centering
\small
\setlength{\tabcolsep}{2pt}
\begin{minipage}[t]{0.555\linewidth}
\vspace{0pt}
{\raggedright\textit{(a) Repairing translation errors}\par}\smallskip

\begin{tabular*}{\linewidth}{@{\extracolsep{\fill}}lrrrr@{}}
\toprule
\textbf{Method} & \textbf{Repaired} & \textbf{Preserved} & \textbf{Tokens} & \textbf{Accepted} \\
\midrule
Direct repair & 0/5 & 5/5 & 3.751 & 5/10 \\
SWE-agent & 0/5 & 5/5 & 12.678 & 5/10 \\
MatchFixAgent & 0/5 & 5/5 & 14.351 & 5/10 \\
LaDiM & \textbf{4/5} & 5/5 & 13.207 & \textbf{9/10} \\
\bottomrule
\end{tabular*}


\end{minipage}\hfill
\begin{minipage}[t]{0.415\linewidth}
\vspace{0pt}
{\raggedright\textit{(b) JAX repair and conversion}\par}\smallskip
\begin{tabular*}{\linewidth}{@{\extracolsep{\fill}}lccc@{}}
\toprule
\textbf{Method} & \textbf{Accepted} & \textbf{Calls} & \textbf{Tokens} \\
\midrule
LaDiM         & 6/6 & 76 & 0.674 \\
Direct repair & 6/6 & 57 & 0.353 \\
Ivy                & 0/6 & 0 & 0 \\
\texttt{torch2jax} & 0/6 & 0 & 0 \\
\bottomrule
\end{tabular*}
\end{minipage}
\end{table}

\begin{figure}[t]
\centering
\includegraphics[width=\linewidth]{figures/tanh-repair-case/tanh_repair_case.pdf}
\caption{Repairing a silent gradient fault. (a) Layered Diagnosis identifies Tanh as the gradient break; the highlighted code adds its missing MindSpore mapping (registration decorators omitted). (b) Cumulative investigation and repair tokens. LaDiM repairs the program, while both baselines remain unrepaired after 40 calls.}
\label{fig:tanh-case}
\end{figure}

\textbf{Case Study: From a Silent Gradient Fault to an Accepted Repair.} \cref{fig:tanh-case} traces a failing translation from \cref{tab:repair-results}a, an autoencoder-style classifier whose encoder applies Tanh between two linear layers. Execution and forward values match the source and the loss differs by only $1.4\times10^{-6}$, yet the first linear layer receives zero gradients. Because the second layer's gradients are correct, the Verifier isolates Tanh, and a targeted test confirms that the supporting library lacks a MindSpore mapping for it, so Tanh silently runs in PyTorch outside MindSpore differentiation. LaDiM registers MindSpore's \texttt{tanh} and passes public verification at call 23, with acceptance confirmed on three seeds and 1.488 million tokens over all 24 calls. SWE-agent and MatchFixAgent also suspect Tanh but never turn the suspicion into a tested patch, exhausting the 40-call limit after 1.606 and 2.341 million tokens. Such faults are invisible to the loss at the step where they occur: gradient and update checks detect injected faults at step~1 in all 12 runs per fault, whereas a loss check at the 0.02 threshold detects 3/12 gradient-scaling and 4/12 update-suppression runs within 50 steps (\cref{sec:detection-measurements}).

\textbf{Generalization Across Frameworks.}
\label{sec:generalization}
LaDiM repairs all six JAX cases spanning execution, forward, and gradient faults in MLP and CNN models (\cref{tab:repair-results}b). Every case passes single-step SGD verification on three seeds after the first submission. Direct repair also accepts all six cases and uses fewer tokens. Ivy and \texttt{torch2jax} accept none of the faulty candidates after native conversion. Both converters pass evaluation on the corresponding MLP and CNN programs without these faults. In these cases, native conversion preserves the supplied execution or numerical faults, and explicit repair restores agreement in training behavior. The same diagnosis and repair procedure thus works beyond MindSpore with JAX differentiation and parameter updates. \cref{sec:jax-protocol} details the converter evaluations. \cref{sec:native-jax-study} reports the separate study of native JAX translations evaluated over consecutive training steps.

\subsection{Ablation Studies}
\label{sec:ablations}

\textbf{Training signals.} We compare cumulative feedback on 16 fixed candidates from a CNN, an image MLP, a Transformer classifier, and a small causal language model. Each model contributes four candidates with execution, forward, gradient, and update faults, respectively. Sources, candidates, tools, and budgets are fixed. Available feedback determines whether repair begins and stops. Each added signal exposes another fault class, increasing final acceptance under complete verification from 4/16 to 16/16 (\cref{tab:signal-ablation}a). Initially undetected faults decrease from 12 to 8 to 4 to 0 across the four settings. Forward values reveal changed computations, gradients reveal differentiation faults, and updates reveal optimizer faults.
\begin{table}[!htb]
\caption{Training signals and repair components. (a) Cumulative signals on 16 faulty candidates, scored with complete verification. Each $+$ adds a signal. (b) Repair history and evidence handoff on ten translations; parentheses count repairs among five failing programs. (c) Repository components on time series; all settings pass 69 checks across three seeds.}
\label{tab:signal-ablation}
\label{tab:cumulative-components}
\centering
\small
\setlength{\tabcolsep}{3pt}
\begin{minipage}[t]{0.495\linewidth}
\vspace{0pt}
{\raggedright\textit{(a) Training signals}\par}\smallskip

\begin{tabular*}{\linewidth}{@{\extracolsep{\fill}}lrr@{}}
\toprule
\textbf{Feedback} & \textbf{Accepted} & \textbf{Tokens} \\
\midrule
Execution and basic checks & 4/16 & 2.089 \\
\quad + Forward values & 8/16 & 2.603 \\
\quad + Gradients & 12/16 & 3.303 \\
\quad + Parameter updates & \textbf{16/16} & 3.419 \\
\bottomrule
\end{tabular*}


\end{minipage}\hfill
\begin{minipage}[t]{0.495\linewidth}
\vspace{0pt}
{\raggedright\textit{(b) Repair components}\par}\smallskip

\begin{tabular*}{\linewidth}{@{\extracolsep{\fill}}lrr@{}}
\toprule
\textbf{Condition} & \textbf{Accepted} & \textbf{Tokens} \\
\midrule
LaDiM & \textbf{9/10 (4/5)} & 15.548 \\
Continuous conversation & 8/10 (3/5) & 14.516 \\
Without repair history & 5/10 (0/5) & 10.852 \\
\bottomrule
\end{tabular*}


\end{minipage}
\par\medskip





\raggedright\textit{(c) Components on the time series repository}\par\smallskip
\begin{tabular*}{\linewidth}{@{\extracolsep{\fill}}>{\raggedright\arraybackslash}p{4.8cm}rrrrr@{}}
\toprule
& \multicolumn{2}{c}{\textbf{Without component}} & \multicolumn{2}{c}{\textbf{With component}} & \\
\cmidrule(lr){2-3}\cmidrule(lr){4-5}
\textbf{Component} & \textbf{Calls} & \textbf{Tokens} & \textbf{Calls} & \textbf{Tokens} & \textbf{Saved} \\
\midrule
Repository context management & 26 & 6.047 & 22 & 1.937 & \textbf{68.0\%} \\
Repository Structural Analysis & 42 & 2.651 & 30 & 2.127 & \textbf{19.8\%} \\
\bottomrule
\end{tabular*}


\par\smallskip
\begin{minipage}{\linewidth}
Tokens are in millions for investigation and repair; (c) also includes initial translation. Calls in (c) count investigation and repair. Saved is the token reduction. Panel (b) uses editing examples and feedback on malformed edits.
\end{minipage}
\end{table}

\textbf{Repair history and independent evidence handoff.} Repair history and independent evidence handoff improve repair on the ten translated programs (\cref{tab:signal-ablation}b). The LaDiM row enables both components, and the other two rows change one component at a time. With repair history, LaDiM repairs four of five failed programs. Removing that history reduces the count to zero. Independent evidence handoff repairs one more program than a shared investigation and repair conversation, increasing final acceptance from 8/10 to 9/10. All three settings preserve the five programs that pass evaluation before repair. Repair history retains earlier attempts and their verification feedback, while independent evidence handoff lets the Repair Agent reassess the Verifier's evidence in its own conversation.

\textbf{Repository context management and Repository Structural Analysis.} On the time series repository, repository context management reduces end-to-end token use by 68.0\% while preserving acceptance on all 69 checks (\cref{tab:signal-ablation}c). Both conditions retain investigation and independent evidence handoff. A separate comparison on the same repository evaluates LaDiM with and without Repository Structural Analysis. Enabling this component reduces end-to-end token use by 19.8\% and LLM calls from 42 to 30, with all 69 checks still passing. Repository Structural Analysis identifies relevant implementations. Repository context management keeps plans, measurements, and inspected code available throughout the investigation. \cref{sec:component-study} gives the complete program component results, and \cref{sec:repository-component-results} reports both repositories' cumulative comparisons and independent ablations of Repository Structural Analysis and Repair Dependency Graph Planning.



\Needspace{6\baselineskip}
\section{Conclusion}
We introduced LaDiM for preserving training behavior during code migration. Layered Diagnosis traces dependencies among execution, forward values, gradients, and parameter updates. Independent evidence handoff and repair history guide corrective edits, while repository context management coordinates work across shared implementations and entry points. Experiments across frameworks, languages, and repositories show that LaDiM preserves training behavior with lower token costs and repairs translation errors that the repair baselines leave unresolved.
\subsection*{AI use statement}
We used generative AI for preliminary literature research, method implementation, result analysis, and manuscript and figure editing. The authors checked the code, results, and citations and take responsibility for the final work.


\bibliographystyle{iclr2027_conference}

\section*{References}
\begin{thebibliography}{39}
\providecommand{\natexlab}[1]{#1}
\providecommand{\url}[1]{\texttt{#1}}
\expandafter\ifx\csname urlstyle\endcsname\relax
  \providecommand{\doi}[1]{doi: #1}\else
  \providecommand{\doi}{doi: \begingroup \urlstyle{rm}\Url}\fi

\bibitem[Abadi et~al.(2016)Abadi, Barham, Chen, Chen, Davis, Dean, Devin,
  Ghemawat, Irving, Isard, Kudlur, Levenberg, Monga, Moore, Murray, Steiner,
  Tucker, Vasudevan, Warden, Wicke, Yu, and Zheng]{abadi2016tensorflow}
Mart{\'i}n Abadi, Paul Barham, Jianmin Chen, Zhifeng Chen, Andy Davis, Jeffrey
  Dean, Matthieu Devin, Sanjay Ghemawat, Geoffrey Irving, Michael Isard,
  Manjunath Kudlur, Josh Levenberg, Rajat Monga, Sherry Moore, Derek~G. Murray,
  Benoit Steiner, Paul Tucker, Vijay Vasudevan, Pete Warden, Martin Wicke, Yuan
  Yu, and Xiaoqiang Zheng.
\newblock {TensorFlow}: A system for large-scale machine learning.
\newblock In \emph{Proceedings of the 12th USENIX Symposium on Operating
  Systems Design and Implementation}, pp.\  265--283, 2016.
\newblock URL
  \url{https://www.usenix.org/conference/osdi16/technical-sessions/presentation/abadi}.

\bibitem[Chen et~al.(2024{\natexlab{a}})Chen, Su, Zuo, Yang, Yuan, Chan, Yu,
  Lu, Hung, Qian, Qin, Cong, Xie, Liu, Sun, and Zhou]{chen2023agentverse}
Weize Chen, Yusheng Su, Jingwei Zuo, Cheng Yang, Chenfei Yuan, Chi-Min Chan,
  Heyang Yu, Yaxi Lu, Yi-Hsin Hung, Chen Qian, Yujia Qin, Xin Cong, Ruobing
  Xie, Zhiyuan Liu, Maosong Sun, and Jie Zhou.
\newblock {AgentVerse}: Facilitating multi-agent collaboration and exploring
  emergent behaviors.
\newblock In \emph{Proceedings of the Twelfth International Conference on
  Learning Representations}, 2024{\natexlab{a}}.
\newblock URL \url{https://openreview.net/forum?id=EHg5GDnyq1}.

\bibitem[Chen et~al.(2024{\natexlab{b}})Chen, Lin, Sch{\"a}rli, and
  Zhou]{chen2023selfdebugging}
Xinyun Chen, Maxwell Lin, Nathanael Sch{\"a}rli, and Denny Zhou.
\newblock Teaching large language models to self-debug.
\newblock In \emph{Proceedings of the Twelfth International Conference on
  Learning Representations}, 2024{\natexlab{b}}.
\newblock URL \url{https://openreview.net/forum?id=KuPixIqPiq}.

\bibitem[Deng et~al.(2023)Deng, Meng, Chen, Liu, Xiang, and
  Chen]{deng2023tensorscope}
Zizhuang Deng, Guozhu Meng, Kai Chen, Tong Liu, Lu~Xiang, and Chunyang Chen.
\newblock Differential testing of cross deep learning framework {APIs}:
  Revealing inconsistencies and vulnerabilities.
\newblock In \emph{32nd USENIX Security Symposium (USENIX Security 23)}, pp.\
  7393--7410. USENIX Association, 2023.
\newblock URL
  \url{https://www.usenix.org/conference/usenixsecurity23/presentation/deng-zizhuang}.

\bibitem[{DJL Contributors}(2026{\natexlab{a}})]{d2ljava2026}
{DJL Contributors}.
\newblock Dive into deep learning with java.
\newblock \url{https://github.com/deepjavalibrary/d2l-java},
  2026{\natexlab{a}}.
\newblock Source repository, accessed September 21, 2026.

\bibitem[{DJL Contributors}(2026{\natexlab{b}})]{djl2026}
{DJL Contributors}.
\newblock Deep java library.
\newblock \url{https://github.com/deepjavalibrary/djl}, 2026{\natexlab{b}}.
\newblock Source repository, accessed September 21, 2026.

\bibitem[{gauravchak}(2026)]{chak2026twotower}
{gauravchak}.
\newblock Two tower models.
\newblock \url{https://github.com/gauravchak/two_tower_models}, 2026.
\newblock Source repository, accessed September 21, 2026.

\bibitem[{Google}(2025)]{google2025torchax}
{Google}.
\newblock {torchax}: A {PyTorch} frontend for {JAX}.
\newblock \url{https://google.github.io/torchax/}, 2025.
\newblock Accessed: 2026-06-02.

\bibitem[Guan et~al.(2025)Guan, Yin, Peng, and Ni]{guan2025repotransagent}
Ziqi Guan, Xin Yin, Zhiyuan Peng, and Chao Ni.
\newblock {RepoTransAgent}: Multi-agent {LLM} framework for repository-aware
  code translation.
\newblock \emph{arXiv preprint arXiv:2508.17720}, 2025.
\newblock \doi{10.48550/arXiv.2508.17720}.
\newblock URL \url{https://arxiv.org/abs/2508.17720}.

\bibitem[He et~al.(2025)He, Chen, Yang, Zhao, Yin, Kang, Lin, Rajmohan, and
  Zhang]{he-etal-2025-execoder}
Minghua He, Yue Chen, Fangkai Yang, Pu~Zhao, Wenjie Yin, Yu~Kang, Qingwei Lin,
  Saravan Rajmohan, and Dongmei Zhang.
\newblock {ExeCoder}: Empowering large language models with executability
  representation for code translation.
\newblock In \emph{Proceedings of the 2025 Conference on Empirical Methods in
  Natural Language Processing}, pp.\  7099--7125, Suzhou, China, 2025.
  Association for Computational Linguistics.
\newblock \doi{10.18653/v1/2025.emnlp-main.362}.
\newblock URL \url{https://aclanthology.org/2025.emnlp-main.362/}.

\bibitem[Hong et~al.(2024)Hong, Zhuge, Chen, Zheng, Cheng, Zhang, Wang, Wang,
  Yau, Lin, Zhou, Ran, Xiao, Wu, and Schmidhuber]{hong2023metagpt}
Sirui Hong, Mingchen Zhuge, Jiaqi Chen, Xiawu Zheng, Yuheng Cheng, Ceyao Zhang,
  Jinlin Wang, Zili Wang, Steven Ka~Shing Yau, Zijuan Lin, Liyang Zhou, Chenyu
  Ran, Lingfeng Xiao, Chenglin Wu, and J{\"u}rgen Schmidhuber.
\newblock {MetaGPT}: Meta programming for a multi-agent collaborative
  framework.
\newblock In \emph{Proceedings of the Twelfth International Conference on
  Learning Representations}, 2024.
\newblock URL \url{https://openreview.net/forum?id=VtmBAGCN7o}.

\bibitem[Ibrahimzada et~al.(2026)Ibrahimzada, Paulsen, Jabbarvand, Dodds, and
  Kroening]{ibrahimzada2025matchfixagent}
Ali~Reza Ibrahimzada, Brandon Paulsen, Reyhaneh Jabbarvand, Joey Dodds, and
  Daniel Kroening.
\newblock {MatchFixAgent}: Language-agnostic autonomous repository-level code
  translation validation and repair.
\newblock In \emph{Forty-third International Conference on Machine Learning},
  2026.
\newblock URL \url{https://openreview.net/forum?id=MuyXpH3GL1}.

\bibitem[{Ivy Contributors}(2026)]{ivy2026transpiler}
{Ivy Contributors}.
\newblock {Ivy}: Convert machine learning code between frameworks.
\newblock \url{https://github.com/unifyai/ivy}, 2026.
\newblock Accessed: 2026-09-22.

\bibitem[{jinglescode}(2026)]{jingles2026timeseries}
{jinglescode}.
\newblock Time series forecasting with pytorch.
\newblock \url{https://github.com/jinglescode/time-series-forecasting-pytorch},
  2026.
\newblock Source repository, accessed September 21, 2026.

\bibitem[Le et~al.(2014)Le, Afshari, and Su]{le2014emi}
Vu~Le, Mehrdad Afshari, and Zhendong Su.
\newblock Compiler validation via equivalence modulo inputs.
\newblock In \emph{Proceedings of the 35th ACM SIGPLAN Conference on
  Programming Language Design and Implementation}, pp.\  216--226. ACM, 2014.
\newblock \doi{10.1145/2594291.2594334}.
\newblock URL \url{https://doi.org/10.1145/2594291.2594334}.

\bibitem[Ma et~al.(2018)Ma, Juefei-Xu, Zhang, Sun, Xue, Li, Chen, Su, Li, Liu,
  Zhao, and Wang]{ma2018deepgauge}
Lei Ma, Felix Juefei-Xu, Fuyuan Zhang, Jiyuan Sun, Minhui Xue, Bo~Li, Chunyang
  Chen, Ting Su, Li~Li, Yang Liu, Jianjun Zhao, and Yadong Wang.
\newblock {DeepGauge}: Multi-granularity testing criteria for deep learning
  systems.
\newblock In \emph{Proceedings of the 33rd ACM/IEEE International Conference on
  Automated Software Engineering}, pp.\  120--131. ACM, 2018.
\newblock \doi{10.1145/3238147.3238202}.
\newblock URL \url{https://doi.org/10.1145/3238147.3238202}.

\bibitem[Macedo et~al.(2025{\natexlab{a}})Macedo, Tian, and
  Adams]{macedo2025codetransengine}
Marcos Macedo, Yuan Tian, and Bram Adams.
\newblock {CodeTransEngine}: Ready-to-use backend for {LLM}-based code
  translation.
\newblock In \emph{ICLR 2025 Third Workshop on Deep Learning for Code},
  2025{\natexlab{a}}.
\newblock URL \url{https://openreview.net/forum?id=9WHJuzjPJf}.

\bibitem[Macedo et~al.(2025{\natexlab{b}})Macedo, Tian, Nie, Cogo, and
  Adams]{macedo2024intertrans}
Marcos Macedo, Yuan Tian, Pengyu Nie, Filipe~R. Cogo, and Bram Adams.
\newblock {InterTrans}: Leveraging transitive intermediate translations to
  enhance {LLM}-based code translation.
\newblock In \emph{Proceedings of the 47th International Conference on Software
  Engineering}, pp.\  2819--2831. IEEE, 2025{\natexlab{b}}.
\newblock \doi{10.1109/ICSE55347.2025.00236}.
\newblock URL \url{https://doi.org/10.1109/ICSE55347.2025.00236}.

\bibitem[{MindSpore}(2020)]{mindspore2020whitepaper}
{MindSpore}.
\newblock {MindSpore}: An all-scenario deep learning computing framework.
\newblock
  \url{https://mindspore-website.obs.cn-north-4.myhuaweicloud.com/white_paper/MindSpore_white_paper_enV1.1.pdf},
  2020.
\newblock White paper, version 1.1.

\bibitem[{MindSpore}(2023)]{mindspore2023mindconverter}
{MindSpore}.
\newblock Migrating from third party frameworks with {MindConverter}.
\newblock
  \url{https://www.mindspore.cn/mindinsight/docs/en/r2.0.0-alpha/migrate_3rd_scripts_mindconverter.html},
  2023.
\newblock Accessed: 2026-06-02.

\bibitem[{MindSpore}(2024)]{mindspore2024gradient}
{MindSpore}.
\newblock Gradient derivation.
\newblock
  \url{https://www.mindspore.cn/docs/en/r2.4.0/migration_guide/model_development/gradient.html},
  2024.
\newblock MindSpore 2.4.0 documentation, Principle Comparison. Accessed:
  2026-09-24.

\bibitem[{MindSpore}(2025)]{mindspore_troubleshooter_loss_compare}
{MindSpore}.
\newblock {TroubleShooter} \texttt{loss\_compare}: Comparing loss curves of two
  training logs.
\newblock
  \url{https://github.com/mindspore-ai/toolkits/blob/e0486eee96a1893b5e1b59f02669804a2d796030/troubleshooter/docs/api/widget/loss_compare.md},
  2025.
\newblock Accessed: 2026-09-25.

\bibitem[{MindSpore}(n.d.)]{mindspore_precision_tuning}
{MindSpore}.
\newblock Large model precision optimization guide.
\newblock
  \url{https://www.mindspore.cn/mindformers/docs/en/r1.7.0/advanced_development/precision_optimization.html},
  n.d.
\newblock MindSpore Transformers r1.7.0 documentation. Accessed: 2026-09-25.

\bibitem[Nguyen et~al.(2013)Nguyen, Qi, Roychoudhury, and
  Chandra]{nguyen2013semfix}
Hoang Duong~Thien Nguyen, Dawei Qi, Abhik Roychoudhury, and Satish Chandra.
\newblock {SemFix}: Program repair via semantic analysis.
\newblock In \emph{Proceedings of the 35th International Conference on Software
  Engineering}, pp.\  772--781. IEEE, 2013.
\newblock \doi{10.1109/ICSE.2013.6606623}.
\newblock URL \url{https://doi.org/10.1109/ICSE.2013.6606623}.

\bibitem[{OpenI}(2025)]{openi2025msadapter}
{OpenI}.
\newblock {MSAdapter}.
\newblock \url{https://openi.pcl.ac.cn/OpenI/MSAdapter/src/tag/v0.6.0}, 2025.
\newblock Version 0.6.0. Accessed: 2026-09-20.

\bibitem[Paszke et~al.(2019)Paszke, Gross, Massa, Lerer, Bradbury, Chanan,
  Killeen, Lin, Gimelshein, Antiga, Desmaison, Kopf, Yang, DeVito, Raison,
  Tejani, Chilamkurthy, Steiner, Fang, Bai, and Chintala]{paszke2019pytorch}
Adam Paszke, Sam Gross, Francisco Massa, Adam Lerer, James Bradbury, Gregory
  Chanan, Trevor Killeen, Zeming Lin, Natalia Gimelshein, Luca Antiga, Alban
  Desmaison, Andreas Kopf, Edward Yang, Zachary DeVito, Martin Raison, Alykhan
  Tejani, Sasank Chilamkurthy, Benoit Steiner, Lu~Fang, Junjie Bai, and Soumith
  Chintala.
\newblock {PyTorch}: An imperative style, high-performance deep learning
  library.
\newblock In \emph{Advances in Neural Information Processing Systems},
  volume~32, pp.\  8024--8035, 2019.
\newblock URL
  \url{https://papers.neurips.cc/paper/9015-pytorch-an-imperative-style-high-performance-deep-learning-library}.

\bibitem[Pei et~al.(2017)Pei, Cao, Yang, and Jana]{pei2017deepxplore}
Kexin Pei, Yinzhi Cao, Junfeng Yang, and Suman Jana.
\newblock {DeepXplore}: Automated whitebox testing of deep learning systems.
\newblock In \emph{Proceedings of the 26th Symposium on Operating Systems
  Principles}, pp.\  1--18. ACM, 2017.
\newblock \doi{10.1145/3132747.3132785}.
\newblock URL \url{https://doi.org/10.1145/3132747.3132785}.

\bibitem[Qian et~al.(2024)Qian, Liu, Liu, Chen, Dang, Li, Yang, Chen, Su, Cong,
  Xu, Li, Liu, and Sun]{qian2023chatdev}
Chen Qian, Wei Liu, Hongzhang Liu, Nuo Chen, Yufan Dang, Jiahao Li, Cheng Yang,
  Weize Chen, Yusheng Su, Xin Cong, Juyuan Xu, Dahai Li, Zhiyuan Liu, and
  Maosong Sun.
\newblock {ChatDev}: Communicative agents for software development.
\newblock In \emph{Proceedings of the 62nd Annual Meeting of the Association
  for Computational Linguistics}, pp.\  15174--15186. Association for
  Computational Linguistics, 2024.
\newblock \doi{10.18653/v1/2024.acl-long.810}.
\newblock URL \url{https://aclanthology.org/2024.acl-long.810/}.

\bibitem[Rozi{\`e}re et~al.(2020)Rozi{\`e}re, Lachaux, Chanussot, and
  Lample]{roziere2020transcoder}
Baptiste Rozi{\`e}re, Marie-Anne Lachaux, Lowik Chanussot, and Guillaume
  Lample.
\newblock Unsupervised translation of programming languages.
\newblock In \emph{Advances in Neural Information Processing Systems},
  volume~33, pp.\  20601--20611, 2020.
\newblock URL
  \url{https://proceedings.neurips.cc/paper/2020/hash/ed23fbf18c2cd35f8c7f8de44f85c08d-Abstract.html}.

\bibitem[Rozi{\`e}re et~al.(2022)Rozi{\`e}re, Zhang, Charton, Harman, Synnaeve,
  and Lample]{roziere2021unittests}
Baptiste Rozi{\`e}re, Jie~M. Zhang, Francois Charton, Mark Harman, Gabriel
  Synnaeve, and Guillaume Lample.
\newblock Leveraging automated unit tests for unsupervised code translation.
\newblock In \emph{Proceedings of the Tenth International Conference on
  Learning Representations}, 2022.
\newblock URL \url{https://openreview.net/forum?id=cmt-6KtR4c4}.

\bibitem[Shinn et~al.(2023)Shinn, Cassano, Berman, Gopinath, Narasimhan, and
  Yao]{shinn2023reflexion}
Noah Shinn, Federico Cassano, Edward Berman, Ashwin Gopinath, Karthik
  Narasimhan, and Shunyu Yao.
\newblock Reflexion: Language agents with verbal reinforcement learning.
\newblock In \emph{Advances in Neural Information Processing Systems}, 2023.
\newblock URL \url{https://arxiv.org/abs/2303.11366}.

\bibitem[Szafraniec et~al.(2023)Szafraniec, Rozi{\`e}re, Leather, Charton,
  Labatut, and Synnaeve]{szafraniec2023compiler}
Marc Szafraniec, Baptiste Rozi{\`e}re, Hugh~James Leather, Francois Charton,
  Patrick Labatut, and Gabriel Synnaeve.
\newblock Code translation with compiler representations.
\newblock In \emph{Proceedings of the Eleventh International Conference on
  Learning Representations}, 2023.
\newblock URL \url{https://openreview.net/forum?id=XomEU3eNeSQ}.

\bibitem[{torch2jax Contributors}(2026)]{torch2jax2026}
{torch2jax Contributors}.
\newblock {torch2jax}: Run {PyTorch} in {JAX}.
\newblock \url{https://github.com/samuela/torch2jax}, 2026.
\newblock Accessed: 2026-09-22.

\bibitem[Wang et~al.(2024)Wang, Chen, Yuan, Zhang, Li, Peng, and
  Ji]{wang2024codeact}
Xingyao Wang, Yangyi Chen, Lifan Yuan, Yizhe Zhang, Yunzhu Li, Hao Peng, and
  Heng Ji.
\newblock Executable code actions elicit better {LLM} agents.
\newblock In \emph{Proceedings of the 41st International Conference on Machine
  Learning}, volume 235 of \emph{Proceedings of Machine Learning Research},
  pp.\  50208--50232, Vienna, Austria, 2024. PMLR.
\newblock URL \url{https://proceedings.mlr.press/v235/wang24h.html}.

\bibitem[Wu et~al.(2024)Wu, Bansal, Zhang, Wu, Li, Zhu, Jiang, Zhang, Zhang,
  Liu, Awadallah, White, Burger, and Wang]{wu2023autogen}
Qingyun Wu, Gagan Bansal, Jieyu Zhang, Yiran Wu, Beibin Li, Erkang Zhu,
  Li~Jiang, Xiaoyun Zhang, Shaokun Zhang, Jiale Liu, Ahmed~Hassan Awadallah,
  Ryen~W. White, Doug Burger, and Chi Wang.
\newblock {AutoGen}: Enabling next-gen {LLM} applications via multi-agent
  conversation.
\newblock In \emph{Proceedings of the First Conference on Language Modeling},
  2024.
\newblock URL \url{https://openreview.net/forum?id=uAjxFFing2}.

\bibitem[Xia et~al.(2023)Xia, Wei, and Zhang]{xia2023aprplm}
Chunqiu~Steven Xia, Yuxiang Wei, and Lingming Zhang.
\newblock Automated program repair in the era of large pre-trained language
  models.
\newblock In \emph{Proceedings of the 45th International Conference on Software
  Engineering}, pp.\  1482--1494. IEEE, 2023.
\newblock \doi{10.1109/ICSE48619.2023.00129}.
\newblock URL \url{https://doi.org/10.1109/ICSE48619.2023.00129}.

\bibitem[Yang et~al.(2023)Yang, Deng, Yao, Tu, Li, and
  Zhang]{yang2023nablafuzz}
Chenyuan Yang, Yinlin Deng, Jiayi Yao, Yuxing Tu, Hanchi Li, and Lingming
  Zhang.
\newblock Fuzzing automatic differentiation in deep-learning libraries.
\newblock In \emph{2023 IEEE/ACM 45th International Conference on Software
  Engineering (ICSE)}, 2023.
\newblock \doi{10.1109/ICSE48619.2023.00105}.
\newblock URL \url{https://arxiv.org/abs/2302.04351}.

\bibitem[Yang et~al.(2024)Yang, Jimenez, Wettig, Lieret, Yao, Narasimhan, and
  Press]{yang2024sweagent}
John Yang, Carlos~E. Jimenez, Alexander Wettig, Kilian Lieret, Shunyu Yao,
  Karthik Narasimhan, and Ofir Press.
\newblock {SWE-agent}: Agent-computer interfaces enable automated software
  engineering.
\newblock \emph{arXiv preprint arXiv:2405.15793}, 2024.
\newblock \doi{10.48550/arXiv.2405.15793}.
\newblock URL \url{https://arxiv.org/abs/2405.15793}.

\bibitem[Yang et~al.(2011)Yang, Chen, Eide, and Regehr]{yang2011csmith}
Xuejun Yang, Yang Chen, Eric Eide, and John Regehr.
\newblock Finding and understanding bugs in {C} compilers.
\newblock In \emph{Proceedings of the 32nd ACM SIGPLAN Conference on
  Programming Language Design and Implementation}, pp.\  283--294. ACM, 2011.
\newblock \doi{10.1145/1993498.1993532}.
\newblock URL \url{https://doi.org/10.1145/1993498.1993532}.

\end{thebibliography}

\clearpage
\appendix
\crefalias{section}{appendix}
\crefalias{subsection}{appendix}
\crefalias{subsubsection}{appendix}

\section{Additional Experimental Protocols}
\label{sec:supplementary-protocols}

\subsection{Task Construction and Accounting}
\label{sec:task-construction}
\begin{table}[!htb]
\caption{Program collections for migration and repair, with their candidate construction, model coverage, and task counts.}
\label{tab:collections}
\centering
\small
\setlength{\tabcolsep}{3pt}

\begin{tabular*}{\linewidth}{@{\extracolsep{\fill}}>{\raggedright\arraybackslash}p{3.1cm}>{\raggedright\arraybackslash}p{3.1cm}>{\raggedright\arraybackslash}p{4.9cm}r@{}}
\toprule
\textbf{Collection} & \textbf{Candidate construction} & \textbf{Model or operation coverage} & \textbf{Tasks} \\
\midrule
PyTorch to MindSpore & Generated translations & Operators, CNNs, MLPs, Transformers, language models & 50 \\
Initial translations & Five fail and five pass their first evaluation & Recurrent, convolutional, attention, and set models & 10 \\
Training signal ablation & Controlled training faults & CNN, image MLP, Transformer, causal language model & 16 \\
JAX repair & Supplied faulty candidates & MLP and CNN & 6 \\
\bottomrule
\end{tabular*}


\end{table}
The main MindSpore collection contains 50 task identifiers. Twelve tasks reuse existing PyTorch source programs, and the remaining tasks use programs authored from public interfaces and declared semantics. The identifiers correspond to 29 distinct combinations of source program and evaluation settings drawn from 24 source files. Conditions with identical source, evaluation settings, initial candidate, method, and inputs are executed once and associated with the corresponding identifiers. Acceptance in the main table uses all 50 identifiers, while token totals count actual calls. LaDiM, SWE-agent, and MatchFixAgent share the same initial translation for each distinct input. Methods may edit the candidate program and create separate temporary tests for investigation. Sources and task specifications remain fixed.

The comparison across languages contains ten DJL application programs and eight textbook examples. LaDiM accepts the fruit classification application using transfer learning and all eight examples. MatchFixAgent accepts the same application and seven examples, while SWE-agent accepts the MNIST LSTM application and seven examples. Test-guided repair accepts seven examples. All repair conditions share the programs produced by initial translation and their generation costs. InterTrans uses its direct Java-to-Python path and an intermediate JavaScript path, with native search and extraction procedures. Its 2.286 million tokens include 0.574 million from incomplete generation calls and 1.712 million from the completed search. The final evaluation compares outputs, losses, gradients, parameter updates, and observable optimizer state over two consecutive training steps.

The common MindSpore initial translations cost 0.429 million tokens. This amount is included once in each repair method's total. Subsequent LaDiM calls use 1.410 million tokens on programs that pass evaluation before repair and 3.320 million on programs that need repair. The corresponding MatchFixAgent totals are 5.136 million and 6.532 million. Across the 29 distinct inputs, the median ratio of LaDiM to MatchFixAgent total tokens is 0.588. LaDiM uses fewer tokens on 19 of the 20 programs that pass before repair and eight of the nine that fail, giving 27/29 overall. LaDiM repairs seven of the nine failing programs at the first submission. Shared initial translation costs cancel when their paired costs are subtracted. MatchFixAgent accepts all tasks at its first submission, which can include multiple analysis and repair calls.

All failed tasks remain in each study's declared denominator, and all failed calls remain in its cost totals. Acceptance at a submission budget scores the current candidate, including any regression after an earlier pass. An edit counts one successfully applied editing operation and can replace several code blocks across files in the declared scope. Investigation and repair share cumulative call, output token, and time budgets across submissions. Token usage is reported in millions and rounded to three decimal places; budgets and per-call limits are given as exact token counts.
\begin{table}[!htb]
\caption{Complete input and output token costs by method for the two program migration collections. Input, output, and total tokens are in millions. Calls and tokens include initial translation and every subsequent call. Acceptance uses the declared task denominator for each collection.}
\label{tab:program-costs}
\centering
\small
\setlength{\tabcolsep}{3pt}

\begin{tabular*}{\linewidth}{@{\extracolsep{\fill}}lrrrrr@{}}
\toprule
\textbf{Method} & \textbf{Calls} & \textbf{Input} & \textbf{Output} & \textbf{Tokens} & \textbf{Accepted} \\
\midrule
\multicolumn{6}{l}{\textit{(a) PyTorch to MindSpore}} \\
Direct LLM & 29 & 0.018 & 0.411 & 0.429 & 29/50 \\
CodeTransEngine~\citep{macedo2025codetransengine} & 29 & 0.022 & 0.257 & 0.280 & 31/50 \\
MSAdapter~\citep{openi2025msadapter} & -- & -- & -- & -- & 15/50 \\
SWE-agent~\citep{yang2024sweagent} & 1,097 & 19.845 & 0.987 & 20.831 & 44/50 \\
MatchFixAgent~\citep{ibrahimzada2025matchfixagent} & 619 & 10.750 & 1.347 & 12.097 & \textbf{50/50} \\
\textbf{LaDiM (ours)} & 326 & 4.298 & 0.861 & 5.159 & \textbf{50/50} \\
\midrule
\multicolumn{6}{l}{\textit{(b) Java/DJL to Python/PyTorch}} \\
Direct LLM & 18 & 0.142 & 0.552 & 0.694 & 1/18 \\
InterTrans~\citep{macedo2024intertrans} & 89 & 0.680 & 1.606 & 2.286 & 1/18 \\
Test-guided repair & 71 & 1.594 & 2.881 & 4.475 & 7/18 \\
SWE-agent & 682 & 23.063 & 1.258 & 24.321 & 8/18 \\
MatchFixAgent & 524 & 26.335 & 2.020 & 28.355 & 8/18 \\
\textbf{LaDiM (ours)} & 345 & 19.470 & 2.009 & 21.479 & \textbf{9/18} \\
\bottomrule
\end{tabular*}


\end{table}

\subsection{Numerical Comparisons and Migration Settings}
\label{sec:migration-settings}
The main MindSpore collection uses absolute/relative tolerances of $10^{-4}/10^{-4}$ for forward values, $10^{-4}/10^{-3}$ for gradients, and $10^{-5}/10^{-3}$ for updates. BF16 matrix multiplication uses absolute tolerance $10^{-2}$. Shapes, data types, initial values, and trainability must match. Parameter and optimizer tasks test their specified enumeration, gradient assignment, or update behavior. Source and target programs begin from corresponding initial states and retain their own state over two consecutive steps for individual migration and three for repositories. The source program and task requirements are available to each method. Acceptance also checks that the target framework executes the required computations.

SWE-agent and MatchFixAgent retain their native algorithms, prompts, tools, parsers, and retry policies. SWE-agent uses its native controller, and MatchFixAgent runs its analysis and repair procedure through the shared DeepSeek tool backend. CodeTransEngine generates its own translation through the native direct path, and MSAdapter uses its native interface mapping.

For each program, each agent has cumulative investigation and repair limits of 40 LLM calls, 120,000 output tokens, 1,800 seconds, and four submissions. Each call allows at most 16,384 output tokens, and LaDiM allows eight calls per investigation or repair stage. Initial translation is separate from this repair budget and allows at most 131,072 output tokens per call. Test-guided repair rewrites the complete candidate, with four repair rounds and at most 131,072 output tokens per call. For each repository, each agent has cumulative limits of 80 calls, 480,000 output tokens, 3,600 seconds, and four submissions, with at most 32,768 output tokens per call. Output budgets include reasoning tokens, and all input usage is recorded.

\subsection{Random Seeds and Initial States}
\label{sec:seeds}

\textbf{Individual migration.} The main MindSpore and Java/DJL comparisons use seed 101 for the evaluation available during repair and seeds 202 and 303 to confirm a passing candidate. We refer to these as the public and confirmation seeds, respectively. The MindSpore evaluator uses NumPy's random generator to sample numerical inputs and training targets. Corresponding parameters are assigned identical initial values generated from a hash of each parameter name; these values remain fixed across the three input seeds. Discrete inputs, masks, and prescribed optimizer-test values retain their task-specific settings. For Java/DJL programs, each paired evaluation supplies the original initial parameters and two preprocessed batches to the candidate for consecutive training steps. The component comparison on the main MindSpore collection uses the same settings.

\textbf{Initial translations and training signals.} The study of ten programs produced by initial translation uses seed 42 for public evaluation and 1042 and 2042 for confirmation. The 16 training-signal candidates use 4101, 5101, and 6101, with 4101 serving as the public seed. The paired evaluator sets the Python, NumPy, and PyTorch random states before constructing the source model and copies its initial state to both executions. A separate CPU PyTorch generator seeded with $s+101$ produces inputs and labels for evaluation seed $s$, including token inputs for text models. The signal study uses all three seeds when checking fault construction and scoring the final candidate with all signals.

\textbf{JAX repair.} The MLP and CNN programs use public seed 6701 and confirmation seeds 6702 and 6703. The evaluator initializes PyTorch's random state with the evaluation seed before constructing the source or candidate model. Inputs and class labels use a separate CPU generator with seed $s+101$.

\textbf{Detection trajectories.} Seeds 300, 301, and 302 define three repeated trajectories for each model. Each seed initializes the model's floating-point parameters, and its initial state is shared across the source and target executions. Image batch order uses seed $s+5000$, text-classification batch order uses $s+6000$, and language-model window selection uses $s+7000$. The fault and synchronization conditions for a given model and seed share the same initial state and 50-step batch sequence.

\textbf{Time series repository.} The main comparison and component ablation use public seed 101 and confirmation seeds 202 and 303. The source model is initialized with the evaluation seed, and its parameters are copied to the target. NumPy generators use seed $s+47$ for the three training batches and seed $s$ for the synthetic price series. The target MindSpore random state also uses $s$. The training script, notebook, and teaching examples are exercised in the public evaluation; confirmation repeats the numerical checks on the remaining two seeds.

\textbf{Recommendation repository.} The public numerical evaluation uses seed 101 for source model initialization and NumPy generation of user identifiers, features, histories, item identifiers, positions, and labels. The source parameters and retrieval corpus are then supplied to the target, whose MindSpore random state uses the same seed. Confirmation with seeds 202 and 303 begins after all public checks pass. The reported recommendation runs do not reach this stage. The original tests use seed 42, including Python, NumPy, and framework random states. For exact comparisons involving the history encoder, the task supplies source parameters initialized with PyTorch seed 42. The one-epoch training command uses seed 101.

\subsection{Training Signal Study}
The training signal and JAX repair studies use the investigation and evidence handoff procedure with editing examples and feedback for correcting edit formats. The signal experiment introduces each of four faults into each of four model classes, producing 16 candidates from four source programs. The execution fault uses an invalid optimizer keyword, the forward fault adds a constant loss offset, the gradient fault scales gradients, and the update fault changes the learning rate.

All feedback conditions retain the initial investigation and current working files across submissions, with the same cumulative call budget. The checks visible to the agent determine whether repair starts and when it stops. Basic checks cover execution, output structure, use of the target framework, preservation of the source program, and the required training records. Final acceptance uses all training signals on three seeds. MindSpore measurements use the \texttt{torch4ms} forward and backward bridge, followed by optimizer updates on the target parameters.

\subsection{JAX Repair and Native Conversion}
\label{sec:jax-protocol}
For JAX, Ivy converts code between frameworks by transpiling the supplied model~\citep{ivy2026transpiler}, while \texttt{torch2jax} maps PyTorch operations and tensors to JAX through native conversion~\citep{torch2jax2026}. We apply each converter once to the supplied faulty candidates and measure whether the existing faults are removed, without LLM repair.

The JAX repair study uses three faulty versions of an MLP and three of a CNN, with an execution, forward, or gradient fault in each model. LaDiM and Direct repair receive the same source programs and faulty candidates. Direct repair keeps investigation and repair in one conversation and permits two LLM calls and one editing operation per repair attempt. Both methods may edit the candidate and create temporary tests, while the source and evaluation settings remain fixed. For each task, both methods allow up to four submissions, 40 LLM calls, 120,000 output tokens, and 1,800 seconds, including investigation. Both methods accept every task at the first submission on seeds 6701, 6702, and 6703. Across the six tasks, LaDiM uses 76 calls and 0.674 million tokens, and Direct repair uses 57 calls and 0.353 million tokens.

TorchAX, a PyTorch frontend for JAX~\citep{google2025torchax}, executes the candidates for both repair methods with JAX arrays. JAX computes gradients, and Optax applies one SGD step. The thresholds are 0.02 for the absolute loss difference, 0.05 for the absolute difference between gradient norms, and 0.03 for the relative parameter update difference. The accepted repairs match the reference initial parameters exactly. The maximum final loss, gradient norm, and relative update differences are $1.192\times10^{-7}$, $1.510\times10^{-8}$, and $1.193\times10^{-6}$.

Ivy produces a JAX/Flax implementation of each supplied model. The \texttt{samuela/torch2jax} implementation converts models and tensors to JAX, followed by JAX differentiation and Optax updates. Each performs a single native conversion of the already faulty candidate. For both converters, the two execution faults retain an extra argument to a linear layer and raise \texttt{TypeError}, leaving numerical measurements n/a. The two forward faults preserve the output scaling error and fail loss, gradient, and update checks. The two gradient faults preserve the interrupted differentiation path: loss passes, while gradients and updates fail. Both converters pass evaluation on the corresponding MLP and CNN programs without these faults. Ivy takes 561.935 seconds for the six faulty candidates and 191.663 seconds for the two programs without faults. The corresponding times for \texttt{torch2jax} are 39.234 and 15.930 seconds. Their LLM calls and tokens are zero.

\subsection{Detection Measurements}
\label{sec:detection-measurements}
The detection study uses four models, three seeds, and 50 training steps. Gradient scaling and partial update suppression are evaluated both with per-step synchronization to the source state and with each execution retaining its own state. With synchronization, each target starts the checked step from the source state. The detection thresholds for the absolute loss difference, absolute difference between gradient norms, and relative parameter update difference are 0.02, 0.05, and 0.03.
\begin{figure}[t]
\centering
\includegraphics[width=0.96\linewidth]{figures/gradient_drift.pdf}
\caption{Detection of (a) gradient scaling and (b) partial update suppression over 50 training steps. Solid lines show the mean over 12 runs (four models $\times$ three seeds), and shaded bands show the minimum and maximum across these runs at each step. Each difference is divided by its detection threshold, so values above the dashed line count as detected. Gradient and update checks detect both faults at step~1 in every run; with the loss threshold of 0.02, the mean loss difference first crosses at steps~31 and~18, respectively.}
\label{fig:gradient-drift}
\end{figure}

\cref{fig:gradient-drift} shows trajectories in which the source and target retain their own states. It reports the mean and range over 12 runs per fault, comprising four models with three seeds each. The mean loss difference crosses its threshold at step 31 for gradient errors and step 18 for update errors. Within 50 steps, the loss difference exceeds its threshold in three of the 12 runs with gradient errors and four of the 12 runs with update errors. Direct gradient or update checks detect the corresponding faults at step~1 in all 12 runs. Translations without these faults pass the recorded checks, with maximum loss difference $3.815\times10^{-6}$.

Loss comparison is an established check in framework migration. MindSpore TroubleShooter's \texttt{loss\_compare} tool plots the loss curves of two training logs and reports their error statistics~\citep{mindspore_troubleshooter_loss_compare}, and the MindSpore Transformers precision guide compares the step-1 loss before gradient norms and updated weights~\citep{mindspore_precision_tuning}. Neither prescribes an automatic threshold, so the loss-only check in \cref{fig:gradient-drift} applies our threshold of 0.02 to the recorded loss differences. Because both faults take effect after the loss is computed, every faulty run has exactly the same loss difference at step~1 as its fault-free counterpart. Lowering the threshold detects the faults earlier (\cref{tab:loss-threshold-sensitivity}): at $10^{-5}$, the loss check detects all 24 faulty runs by step~8 without false alarms, since the largest loss difference among fault-free runs is $3.815\times10^{-6}$. All thresholds are applied post hoc to the same saved trajectories.

\begin{table}[!htb]
\caption{Loss-only detection at different thresholds, computed post hoc on the saved trajectories (12 runs per fault, 50 steps). First detection gives the range of first-crossing steps over detected runs; mean crossing is the first step at which the mean loss difference crosses the threshold; false alarms count fault-free runs that cross the threshold.}
\label{tab:loss-threshold-sensitivity}
\centering
\small
\setlength{\tabcolsep}{3.5pt}
\begin{tabular}{@{}lcccc@{\hspace{10pt}}cccc@{}}
\toprule
 & \multicolumn{4}{c}{(a) Gradient scaling} & \multicolumn{4}{c}{(b) Partial update suppression} \\
\cmidrule(lr){2-5}\cmidrule(l){6-9}
Threshold & Detected & \hdr{First}{detection} & \hdr{Mean}{crossing} & \hdr{False}{alarms} & Detected & \hdr{First}{detection} & \hdr{Mean}{crossing} & \hdr{False}{alarms} \\
\midrule
0.02      & 3/12  & 10--11 & 31 & 0/12 & 4/12  & 6--20 & 18 & 0/12 \\
0.005     & 12/12 & 3--37  & 9  & 0/12 & 8/12  & 3--42 & 6  & 0/12 \\
0.001     & 12/12 & 2--8   & 3  & 0/12 & 9/12  & 2--14 & 2  & 0/12 \\
$10^{-5}$ & 12/12 & 2      & 2  & 0/12 & 12/12 & 2--8  & 2  & 0/12 \\
\bottomrule
\end{tabular}
\end{table}

The detection study uses the following scalar differences:
\begin{align}
\Delta_{\mathrm{loss}} &= |L(P)-L(\hat P)|, \\
\Delta_{\mathrm{grad}} &= \left|\|\nabla_{\theta}L(P)\|_2-\|\nabla_{\hat\theta}L(\hat P)\|_2\right|, \\
\Delta_{\mathrm{param}} &= \frac{\|\Delta\theta-\Delta\hat\theta\|_2}{\|\Delta\theta\|_2+\epsilon}.
\end{align}
Here $\Delta\theta$ and $\Delta\hat\theta$ denote source and target parameter updates, and $\epsilon$ prevents division by zero. \cref{fig:gradient-drift} divides each difference by its detection threshold. The main migration comparison uses the task-specific tensor comparisons defined above.

\section{Evidence Handoff and Agent Integration}
\label{sec:component-study}
The component study uses the same 50 MindSpore task identifiers, initial translations, acceptance checks, and budgets as the main comparison. The continuous conversation condition keeps investigation and repair in one conversation, including earlier messages and reasoning history. Independent evidence handoff supplies the Verifier's evidence to a separate repair conversation. Without repair history, the context for each repair attempt after the first starts from the initial investigation and latest observation. The current files and cumulative budget are retained. Disabling progress reminders leaves the call and submission limits unchanged.
\begin{table}[!htb]
\caption{Component variants and added investigation for baseline agents on the 50 MindSpore task identifiers. Tokens are in millions and include initial translation and all subsequent calls.}
\label{tab:components}
\centering
\small
\setlength{\tabcolsep}{4pt}

\begin{tabular*}{\linewidth}{@{\extracolsep{\fill}}lrr@{}}
\toprule
\textbf{Setting} & \textbf{Accepted} & \textbf{Tokens} \\
\midrule
LaDiM & 50/50 & 5.159 \\
Continuous investigation and repair conversation & 49/50 & 5.603 \\
Without repair history & 50/50 & 4.566 \\
Without progress reminders & 50/50 & 4.562 \\
MatchFixAgent & 50/50 & 12.097 \\
MatchFixAgent with independent investigation & 50/50 & 19.480 \\
SWE-agent & 44/50$^{\dagger}$ & 20.831 \\
SWE-agent with independent investigation & 43/50 & 36.547 \\
\bottomrule
\end{tabular*}


\par\vspace{3pt}
\begin{minipage}{\linewidth}
$\dagger$ SWE-agent without added investigation has one interrupted task with unavailable final acceptance. It remains in the denominator and cost totals. Tasks with identical inputs are executed once as described in \cref{sec:task-construction}.
\end{minipage}
\end{table}
LaDiM accepts 50/50 task identifiers, compared with 49/50 when investigation and repair share one conversation. The variants without repair history and without progress reminders each accept 50/50 while using fewer tokens than LaDiM. Among the 29 distinct inputs represented by these identifiers, 20 translated programs pass evaluation before repair. Only two of the nine failing programs require a second repair submission under LaDiM. Most inputs are therefore accepted before a later repair attempt can use earlier repair history.

The second component study uses the same ten programs produced by initial translation as \cref{tab:repair-results}a, of which five fail evaluation and five pass before repair. Methods may edit these programs and their supporting libraries. LaDiM in \cref{tab:natural-components} uses independent evidence handoff and repair history, with examples of the edit format and corrective feedback for malformed edits. Each other condition changes one component while retaining the candidates, workspace access, acceptance checks, and repair budgets. Removing the edit examples and format correction gives the LaDiM configuration reported in \cref{tab:repair-results}a. All conditions also pass evaluation on the five programs that pass before repair.
\begin{table}[!htb]
\caption{Complete component results on ten programs produced by initial translation. Repaired counts failed programs corrected, and Preserved counts passing programs that remain accepted. Tokens are in millions and cover investigation and repair on all ten programs, including failed attempts.}
\label{tab:natural-components}
\centering
\small
\setlength{\tabcolsep}{3pt}

\begin{tabular*}{\linewidth}{@{\extracolsep{\fill}}>{\raggedright\arraybackslash}p{4.6cm}rrrr@{}}
\toprule
\textbf{Condition} & \textbf{Repaired} & \textbf{Preserved} & \textbf{Tokens} & \textbf{Accepted} \\
\midrule
LaDiM & 4/5 & 5/5 & 15.548 & 9/10 \\
Continuous conversation & 3/5 & 5/5 & 14.516 & 8/10 \\
Without repair history & 0/5 & 5/5 & 10.852 & 5/10 \\
Without progress reminders & 2/5 & 5/5 & 17.325 & 7/10 \\
Without editing format assistance & 4/5 & 5/5 & 13.207 & 9/10 \\
\bottomrule
\end{tabular*}


\end{table}
For SWE-agent and MatchFixAgent, we add the same investigation before repair, with code editing disabled during investigation. The investigation counts toward each baseline agent's total call, token, and time limits. Each baseline agent retains its own history and repair policy.

MatchFixAgent with investigation accepts 50/50 at a higher token total. SWE-agent accepts 43/50 with the added investigation and 44/50 without it. Among the nine distinct programs that fail evaluation before repair, SWE-agent repairs seven without the added investigation and eight with it. With the added investigation, four of the 20 programs that pass before repair fail the final evaluation. Final acceptance in \cref{tab:components} counts both repaired programs and programs that continue to pass after repair.

\section{Additional Repair Collections}
\label{sec:additional-repair}
The controlled MindSpore repair collection contains 50 injected faults in supplied target programs, with task-specific acceptance checks. Direct repair uses the same inspection, testing, and editing tools as LaDiM, while SWE-agent uses its native tools. LaDiM accepts 45/50 within four submissions, compared with 29/50 for SWE-agent and 35/50 for Direct repair. LaDiM's cumulative accepted counts are 29, 41, and 45 with submission limits of one, two, and four, respectively. Total token use is 39.983, 42.579, and 12.459 million for LaDiM, SWE-agent, and Direct repair, respectively.

A separate collection supplies twelve source programs together with faulty target programs, providing the program pairs required by MatchFixAgent. LaDiM accepts 12/12 and MatchFixAgent accepts 11/12, using 6.643 and 14.506 million tokens, respectively. Both receive the complete source program, faulty target program, task specification, and target library. Each method has up to four submissions, 40 LLM calls, 120,000 output tokens, and 1,800 seconds per task. Both repair collections permit edits to the candidate and its supporting library. Their original tests check execution and computed values, without an additional confirmation stage across random seeds.

\section{Natural Translation Repair Protocol}
\label{sec:natural-protocol}
Of the ten programs produced by initial translation, five fail evaluation and five pass before repair. Each method receives the same source, candidate, task requirements, and supporting library. The candidate and library may be edited. The source and acceptance criteria remain fixed.

Target forward computation and differentiation execute in MindSpore through \texttt{torch4ms}. PyTorch then applies optimizer updates to the target parameters. The comparison uses matching inputs and initial parameters, with acceptance confirmed on three seeds. The thresholds are 0.02 for the absolute loss difference, 0.05 for the absolute difference between gradient norms, and 0.03 for the relative parameter update difference. The L2 difference between gradient vectors has an additional threshold of 0.05, alongside the specified layer and parameter checks.

The LLM call counts in \cref{tab:repair-results}a cover all ten programs, including failed repairs and investigation of programs that pass before repair. LaDiM uses 193 calls, Direct repair uses 120, SWE-agent uses 386, and MatchFixAgent uses 265. LaDiM uses 13.207 million tokens.

In the LSTM example in \cref{sec:fault-characterization}, LaDiM repairs two errors in the supporting library. The first submission leaves an unregistered recurrent operator unresolved. LaDiM reproduces the failure, implements the recurrence in MindSpore, and registers the recurrent operator. The second submission exposes recursive tensor indexing. A focused test identifies the repeated calls in tensor indexing, which the third submission corrects. On the public evaluation seed, the final absolute loss difference is $2.38\times10^{-7}$, gradient vector L2 difference is $4.19\times10^{-7}$, and relative parameter update difference is $8.11\times10^{-7}$. The example uses 32 LLM calls.

In the autoencoder example of \cref{fig:tanh-case}, the supporting library registers the low-level Tanh operator but not \texttt{torch.tanh} or \texttt{torch.nn.functional.tanh}. Unregistered calls fall back to PyTorch on copied tensors, which preserves forward values but disconnects the MindSpore differentiation graph. Before repair, on the public evaluation seed, layer outputs differ by at most $1.5\times10^{-8}$ and the absolute loss difference is $1.43\times10^{-6}$. The gradients of the first linear layer's weight and bias are zero, against norms of 0.342 and 0.109 in the source, while the other six parameter gradients differ by at most $6.2\times10^{-7}$. The gradient vector L2 difference is 0.359 and the relative parameter update difference is 0.771, whereas the absolute gradient-norm difference, 0.046, stays below its threshold. LaDiM's first submission contains no edits and fails. Its diagnostic test compares summed absolute input gradients under MindSpore differentiation: 4.894 for native \texttt{tanh} and zero for dispatched Tanh, while dispatched ReLU, sigmoid, and multiplication all pass nonzero gradients. The patch adds a four-line registration to the library, and the second submission is accepted on all three seeds, with gradient vector L2 difference $1.01\times10^{-6}$ and relative parameter update difference $6.42\times10^{-6}$ on the public seed. SWE-agent and MatchFixAgent both attribute the zero gradients to Tanh, and MatchFixAgent even drafts the same registration, yet neither applies an edit to the library within 40 calls, so both leave the program unrepaired. This case illustrates the completion condition in \cref{sec:repair-agent}: the Repair Agent must apply an edit and test after that edit before completing the repair stage.

\section{Repository Protocols and Detailed Results}
\label{sec:repository-results}
All methods receive identical copies of the source repositories and the same initial translations, task requirements, and acceptance procedures. The time series translation has three generation calls totaling 0.164 million tokens. The recommendation translation uses 20 calls totaling 0.368 million tokens. Translation follows the source repository's imports, and its cost includes failed or truncated attempts. In \cref{tab:repository-costs}, each method's total includes the common initial translation once and all subsequent investigation and repair calls.
\begin{table}[!htb]
\caption{Training checks and complete costs for repository migration. Each training signal has nine checks for time series, from three steps on three seeds, and 18 for recommendation, from six model execution paths and three steps on the public seed. Calls, input tokens, and output tokens cover investigation and repair. Tokens reports the total including the common initial translation. Tokens are in millions.}
\label{tab:repository-costs}
\centering
\small
\setlength{\tabcolsep}{2pt}

\begin{tabular*}{\linewidth}{@{\extracolsep{\fill}}lrrrrrrr@{}}
\toprule
\textbf{Method} & \textbf{Gradients} & \thead{Parameter\\updates} & \textbf{Calls} & \textbf{Submissions} & \textbf{Input} & \textbf{Output} & \textbf{Tokens} \\
\midrule
\multicolumn{8}{l}{\textit{Time series}} \\
LaDiM & 9/9 & 9/9 & 35 & 2 & 2.345 & 0.105 & 2.614 \\
SWE-agent & 9/9 & 9/9 & 80 & 1 & 7.985 & 0.051 & 8.200 \\
MatchFixAgent & 9/9 & 9/9 & 69 & 2 & 13.306 & 0.128 & 13.598 \\
\midrule
\multicolumn{8}{l}{\textit{Recommendation}} \\
LaDiM & 18/18 & 15/18 & 80 & 4 & 8.148 & 0.295 & 8.811 \\
SWE-agent & 17/18 & 12/18 & 80 & 1 & 7.890 & 0.058 & 8.317 \\
MatchFixAgent & 18/18 & 16/18 & 80 & 1 & 8.215 & 0.095 & 8.679 \\
\bottomrule
\end{tabular*}


\end{table}
The time series protocol compares the source and translated core models for three consecutive Adam steps from corresponding parameters, with dropout disabled and input shape $4\times20\times1$. Preprocessing and prediction absolute tolerances are $10^{-9}$ and $10^{-4}$. The tolerances are 0.02 for the absolute loss difference, 0.05 for the L2 difference between gradient vectors, and 0.03 for the relative parameter update difference. Learning rate behavior is checked around its schedule boundaries. The public evaluation executes the 100-epoch training workflow, notebook cells, and teaching snippets in their source order. Each entry point must produce the required predictions and plots and compute gradients through the target framework's native differentiation. The notebook and snippets test execution, while the core model provides the numerical comparison of source and target training. All three methods pass 23 behavior checks on each of the public and two confirmation seeds, giving 69/69 overall.

The recommendation protocol compares six model execution paths in the source and translated programs over three sequential Adam updates. It also checks whether the translated program reproduces source exceptions and empty returns. Both the numerical evaluation and training command use the candidate's function for constructing the optimizer, with learning rate 0.001. The source Adam settings use decay coefficients 0.9 and 0.999 and epsilon $10^{-8}$. Representation absolute tolerance is $10^{-4}$. Loss comparisons use absolute tolerance 0.002 and relative tolerance $10^{-4}$; gradient comparisons use $10^{-4}$ and 0.005, respectively. Relative parameter update tolerance is 0.03. The evaluation includes all ten original tests and one epoch of the training command on 16 synthetic examples.

\cref{tab:repository-checks} reports the recommendation checks on the public seed. The 114 numerical comparisons comprise 108 checks over six paths, three steps, and six training signals, plus six inference retrieval checks. Forward loss is returned by the model's forward computation; training return loss is returned by the training function on the same batch. Each has 18 checks. Twenty-four execution, step count, and native differentiation checks and seven source behavior checks bring the behavior check total to 145.

SWE-agent and MatchFixAgent each lack one retrieval measurement after reward model inference fails. These values are n/a, and the missing checks remain in the retrieval, numerical, and behavior denominators of six, 114, and 145. MatchFixAgent has no confirmed passes among the ten original tests because an unmigrated PyTorch import in its history encoder test prevents test collection under the native MindSpore requirement. This is reported as 0/10. No method passes all public checks, so the two confirmation seeds are not evaluated on this repository.

File coverage checks that the translated repository contains the corresponding Python files and public function and class names, preserves the original test function names, and has valid Python syntax in files and README code blocks. It also checks use of MindSpore in place of the source or alternate frameworks. Documentation and dependency checks require the README and dependency list to identify MindSpore, and the dependency list to exclude the source and alternate frameworks. These repository checks are reported separately from the 145 behavior checks.

LaDiM passes all 18 forward loss, 18 training return loss, and 18 gradient checks, together with 15/18 parameter update checks, in the main recommendation comparison. Item representations exceed the comparison tolerance in all six model execution paths at steps two and three. LaDiM also completes reward model inference, retrieval, the training command, and all ten original tests.
\begin{table}[!t]
\caption{Recommendation repository checks passed on the public seed. The 114 numerical checks are included in the 145 behavior checks. The six training signal rows divide the 108 training checks, with 18 per signal. One unmeasured retrieval check remains in the denominators for SWE-agent and MatchFixAgent.}
\label{tab:repository-checks}
\centering
\small
\setlength{\tabcolsep}{4pt}

\begin{tabular*}{\linewidth}{@{\extracolsep{\fill}}lrrr@{}}
\toprule
\textbf{Check} & \textbf{LaDiM} & \textbf{SWE-agent} & \textbf{MatchFixAgent} \\
\midrule
Behavior checks passed & 130/145 & 104/145 & 129/145 \\
Numerical checks & 99/114 & 74/114 & 99/114 \\
Training checks & 93/108 & 69/108 & 94/108 \\
Inference retrieval & 6/6 & 5/6 & 5/6 \\
\midrule
User representations & 18/18 & 16/18 & 18/18 \\
Item representations & 6/18 & 6/18 & 6/18 \\
Forward loss & 18/18 & 18/18 & 18/18 \\
Training return loss & 18/18 & 0/18 & 18/18 \\
Gradients & 18/18 & 17/18 & 18/18 \\
Parameter updates & 15/18 & 12/18 & 16/18 \\
\midrule
Original tests & 10/10 & 10/10 & 0/10 \\
Training command & 1/1 & 1/1 & 1/1 \\
File coverage & 1/1 & 1/1 & 0/1 \\
Documentation and dependencies & 1/1 & 1/1 & 0/1 \\
\bottomrule
\end{tabular*}


\end{table}

\section{Repository Context and Evidence}
\label{sec:repository-context}
Repository context management tracks changes to the dependency plan and candidate files separately. Every recorded test identifies the code version on which it ran. After a code change, measurements from the old version no longer describe the current candidate. A new result replaces the previous result for the same test, including when the new test fails. The validity of each work unit's checks also depends on its files, selected tests, and prerequisite units. Changes to these files or dependencies invalidate the affected checks and those of declared dependent units. A unit checked without a selected test records syntax validation only. Final acceptance requires evaluation of the complete repository.

Full tool observations, previously read code, and conversations are stored outside the editable workspace and remain available for later retrieval. When rebuilding the conversation context, LaDiM includes the current dependency plan and translation progress, the complete structured findings, hypotheses and next steps, the latest measurements, and the most recent tool results. It also includes up to 120,000 characters of previously read code that still matches the current files and is relevant to the active work unit. Context is rebuilt for independent evidence handoff, when switching between work units, or when messages exceed 75\% of the configured context capacity. Continuing a stage or selecting the first work unit does not itself trigger reconstruction.

For the repository experiments, investigation allows up to six LLM calls and 64,000 output tokens. The remaining calls and output tokens are allocated over the remaining permitted repair submissions, with at most 16 calls in the first repair stage. Investigation and all repair stages share the cumulative limits for the whole repository.

\section{Repository Component Results}
\label{sec:repository-component-results}

\subsection{Cumulative Components}
We add Verifier investigation, independent evidence handoff, and repository context management cumulatively to the Repair Agent on the time series and recommendation repositories (\cref{tab:repository-cumulative-full}). Each condition runs once with the same source, initial candidate, evaluation inputs, seeds, thresholds, full training signals, and total budgets. The base condition retains the repair prompts, ordinary inspection and editing tools, syntax checks, and progress reminders. After adding investigation, investigation and repair share one conversation until independent evidence handoff is also added. All four conditions accept the time series repository with 69/69 checks and pass 131/145 checks on recommendation. Each condition therefore fully accepts one of the two repositories.
\begin{table}[!htb]
\caption{Cumulative agent components on both repositories. Checks reports passed behavior checks. Calls cover investigation and repair, and tokens also include the shared initial translation. Each row adds a component to the preceding row. Tokens are in millions.}
\label{tab:repository-cumulative-full}
\centering
\small
\setlength{\tabcolsep}{3pt}



\begin{tabular*}{\linewidth}{@{\extracolsep{\fill}}>{\raggedright\arraybackslash}p{5.4cm}rrrrrr@{}}
\toprule
& \multicolumn{3}{c}{\textbf{Time series}} & \multicolumn{3}{c}{\textbf{Recommendation}} \\
\cmidrule(lr){2-4}\cmidrule(l){5-7}
\textbf{Condition} & \textbf{Checks} & \textbf{Calls} & \textbf{Tokens} & \textbf{Checks} & \textbf{Calls} & \textbf{Tokens} \\
\midrule
Repair Agent & 69/69 & 26 & 6.128 & 131/145 & 62 & 8.336 \\
\quad + Verifier investigation & 69/69 & 18 & 3.896 & 131/145 & 59 & 8.913 \\
\quad + Independent evidence handoff & 69/69 & 26 & 6.047 & 131/145 & 68 & 9.556 \\
\quad + Repository context management & 69/69 & 22 & 1.937 & 131/145 & 80 & 10.191 \\
\bottomrule
\end{tabular*}


\end{table}
On time series, LaDiM with all three added components uses 1.937 million total tokens, 68.4\% fewer than the base Repair Agent. Adding repository context management to investigation and independent evidence handoff reduces tokens from 6.047 million to 1.937 million. Investigation and repair input tokens fall from 5.731 million to 1.700 million. The two truncated LLM responses in the condition without repository context management count toward its call and token totals.

On recommendation, adding the components increases token use while all four conditions pass 131/145 checks and all ten original tests. The 14 unresolved checks comprise 12 item representation checks over six model execution paths at steps two and three, and two parameter update checks. The first three conditions in the table reach the conversation context limit after 62, 59, and 68 calls, respectively. LaDiM with all components rebuilds the conversation context six times and continues to the 80-call limit.

\subsection{Repository Structural Analysis}
The comparison with and without Repository Structural Analysis retains Repair Dependency Graph Planning, notebook tools, ordinary file inspection, independent evidence handoff, and evidence management in both conditions. Removing Repository Structural Analysis removes the inventory of files, imports, symbols, and notebook cells, together with the list of files not yet assigned to a work unit. With the component enabled, the agent requests structural analysis when needed. Initialization and context reconstruction do not perform this analysis. Each condition runs once on each repository with the same initial translation and repository budgets.
\begin{table}[!htb]
\caption{Repository comparisons that remove one component at a time. Both comparisons use the same LaDiM run with all components for each repository. Checks reports passed behavior checks. Calls cover investigation and repair, and tokens include initial translation. Tokens are in millions.}
\label{tab:repository-independent}
\centering
\small
\setlength{\tabcolsep}{3pt}





\begin{tabular*}{\linewidth}{@{\extracolsep{\fill}}>{\raggedright\arraybackslash}p{5.4cm}rrrrrr@{}}
\toprule
& \multicolumn{3}{c}{\textbf{Time series}} & \multicolumn{3}{c}{\textbf{Recommendation}} \\
\cmidrule(lr){2-4}\cmidrule(l){5-7}
\textbf{Condition} & \textbf{Checks} & \textbf{Calls} & \textbf{Tokens} & \textbf{Checks} & \textbf{Calls} & \textbf{Tokens} \\
\midrule
LaDiM & 69/69 & 30 & 2.127 & 131/145 & 80 & 9.962 \\
Without Repository Structural Analysis & 69/69 & 42 & 2.651 & 131/145 & 80 & 7.172 \\
Without Repair Dependency Graph Planning & 69/69 & 21 & 1.804 & 130/145 & 69 & 9.214 \\
\bottomrule
\end{tabular*}


\end{table}
Repository Structural Analysis reduces time series calls from 42 to 30 and total tokens by 19.8\%, with both conditions passing all 69 checks after two submissions. On recommendation, both pass 131/145 checks. With structural analysis, this result is reached at the third submission instead of the fourth. Enabling structural analysis uses 9.962 million total tokens, 38.9\% more than the 7.172 million used without it. The two conditions have the same final 14 failures described above. The runs with all components in this study and the cumulative study were executed separately, so each table reports the costs of its own comparison.

\subsection{Repair Dependency Graph Planning}
Removing Repair Dependency Graph Planning removes the work-unit plan, checks that prerequisite units are ready, restrictions on which files belong to each unit, invalidation of checks through dependencies, and context reconstruction when the active unit changes. The condition retains Repository Structural Analysis, ordinary file and notebook tools, code editing and syntax checks, independent evidence handoff, evidence associated with code versions, and reconstruction at the context limit. The entire permitted target repository remains editable. The conditions with and without planning in \cref{tab:repository-independent} use the same source, initial candidate, inputs, seeds, thresholds, underlying method implementation, and budgets. Each condition runs once on each repository. The run with all components is shared with the structural analysis comparison.

Without Repair Dependency Graph Planning, time series completes all 69 checks after one submission, using 21 calls and 1.804 million total tokens. This is 15.2\% fewer tokens than LaDiM with all components. On recommendation, removing planning uses 7.5\% fewer tokens but passes 130/145 checks, one fewer than LaDiM with all components. The additional failure is the parameter update at step three in the model with position debiasing. Both conditions fully accept one of the two repositories.

Without planning, submissions one through four pass 128, 103, 133, and 130 of the 145 recommendation checks. A change to a shared training entry point before the second submission causes the first regression. The final submission again reduces the number of passed checks from 133 to 130. With all components, the corresponding trajectory is n/a, 122, 131, and 131. The first submission reports eight execution failures separately from the 145 protocol checks, leaving all protocol measurements unavailable. The first submission without planning and the second submission with all components each lack one retrieval measurement; both final candidates have all 145 measurements. Removing planning exhausts the budget of 480,000 output tokens after 69 calls, while LaDiM with all components uses all 80 calls. Final scores use the last submitted candidate, including these regressions.

\section{Native JAX Migration with Consecutive Training Steps}
\label{sec:native-jax-study}
This study uses twelve public model workloads from PyTorch examples, minGPT, and torchvision.\footnote{\url{https://github.com/pytorch/examples}, \url{https://github.com/karpathy/minGPT}, and \url{https://github.com/pytorch/vision}.} They cover image classification, super-resolution, a variational autoencoder, image generation, recurrent and Transformer language models, time series, two policy learning models, graph attention, GPT, and ResNet. Initial translation returns code for ten workloads. Eight translated programs pass evaluation on all three seeds before repair, and two fail. Initial translation of the image generation and Transformer language modeling workloads exhausts the output budget without returning code. All four methods receive the same two failing translations, for the time series model and ResNet, for repair.
\begin{table}[!htb]
\caption{Repair of two failing native JAX translations from twelve source workloads. Calls and Tokens cover repair. Two tasks and All sources report total tokens including 0.031 and 0.140 million shared initial translation tokens, respectively, counted once. Each method accepts 1/2 repair tasks and 9/12 source workloads overall. Tokens are in millions.}
\label{tab:native-jax}
\centering
\small
\setlength{\tabcolsep}{3pt}



\begin{tabular*}{\linewidth}{@{\extracolsep{\fill}}lrrrr@{}}
\toprule
\textbf{Method} & \textbf{Calls} & \textbf{Tokens} & \textbf{Two tasks} & \textbf{All sources} \\
\midrule
LaDiM & 35 & 1.840 & 1.871 & 1.980 \\
Direct repair & 28 & 0.761 & 0.792 & 0.901 \\
MatchFixAgent & 23 & 0.259 & 0.290 & 0.399 \\
SWE-agent & 66 & 1.014 & 1.045 & 1.153 \\
\bottomrule
\end{tabular*}


\end{table}
The translated programs implement the forward computation in native JAX and explicitly represent model buffers, which store state other than trainable parameters. The evaluator independently invokes JAX differentiation and the specified Optax optimizer for three consecutive steps on seeds 8101, 8102, and 8103. Source parameters, buffers, and batches are supplied to each method. The source and JAX programs start from exactly matching states. The evaluator compares their full outputs, losses, gradients, parameter updates, parameters, and buffers elementwise in float64, with absolute tolerance $10^{-7}$ and relative tolerance $10^{-5}$. Tensor names, shapes, and data types must also match. Each method on each repair task has a cumulative budget of 40 LLM calls, four submissions, 120,000 output tokens, and 1,800 seconds, including investigation. Each initial translation has one call with at most 16,384 output tokens.

All four methods repair shape errors in the time series candidate's input slicing and output concatenation, after which the programs pass all three steps on all three seeds. LaDiM, MatchFixAgent, and SWE-agent pass at their first submission, and Direct repair passes at its second. LaDiM, Direct repair, and SWE-agent also turn the truncated ResNet candidate into executable code. These three ResNet programs pass all first-step numerical checks. Later steps contain gradient and parameter update differences, followed by output and buffer differences. MatchFixAgent stops when parsing the truncated ResNet candidate encounters a syntax error. Its numerical results are n/a, and its four completed calls remain in the costs.

The eight programs that pass before repair are CNN classification, super-resolution, the variational autoencoder, the recurrent language model, both policy learning models, graph attention, and GPT. Adding the repaired time series model gives 9/12 accepted source workloads for every method. Initial translation of all twelve sources uses 0.029 million input and 0.110 million output tokens, totaling 0.140 million, including both generation failures. The ResNet results show why consecutive training checks are necessary even when the first step agrees.


\end{document}
